% Ngân hàng đề Toán 12 (KNTT) -- Chương II, Bài 7, Dạng 06
% Dạng: Tọa độ trung điểm, trọng tâm, điểm chia đoạn thẳng
%% Phương án ĐẦU TIÊN trong \cauTN luôn là đáp án đúng (gói trondethi tự đảo).

\cauTN[id=C2-B7-D06-NB-001,muc=NB,loigiai={Trung điểm $I$ của $AB$ có tọa độ $I\left(\dfrac{x_A+x_B}{2};\dfrac{y_A+y_B}{2};\dfrac{z_A+z_B}{2}\right)=(2;-\dfrac{3}{2};\dfrac{1}{2})$.}]{Trong không gian $Oxyz$, cho $A(2;-3;-1)$, $B(2;0;2)$. Tọa độ trung điểm $I$ của đoạn thẳng $AB$ là}
  {$(2;-\dfrac{3}{2};\dfrac{1}{2})$.}
  {$(2;-\dfrac{1}{2};\dfrac{1}{2})$.}{$(1;-\dfrac{3}{2};\dfrac{1}{2})$.}{$(2;-\dfrac{5}{2};\dfrac{1}{2})$.}

\cauTN[id=C2-B7-D06-NB-002,muc=NB,loigiai={Trung điểm $I$ của $AB$ có tọa độ $I\left(\dfrac{x_A+x_B}{2};\dfrac{y_A+y_B}{2};\dfrac{z_A+z_B}{2}\right)=(-3;-1;1)$.}]{Trong không gian $Oxyz$, cho $A(-3;-1;-1)$, $B(-3;-1;3)$. Tọa độ trung điểm $I$ của đoạn thẳng $AB$ là}
  {$(-3;-1;1)$.}
  {$(-3;-1;0)$.}{$(-3;0;1)$.}{$(-4;-1;1)$.}

\cauTN[id=C2-B7-D06-NB-003,muc=NB,loigiai={Trung điểm $I$ của $AB$ có tọa độ $I\left(\dfrac{x_A+x_B}{2};\dfrac{y_A+y_B}{2};\dfrac{z_A+z_B}{2}\right)=(-\dfrac{1}{2};-1;\dfrac{1}{2})$.}]{Trong không gian $Oxyz$, cho $A(-3;-1;1)$, $B(2;-1;0)$. Tọa độ trung điểm $I$ của đoạn thẳng $AB$ là}
  {$(-\dfrac{1}{2};-1;\dfrac{1}{2})$.}
  {$(\dfrac{1}{2};-1;\dfrac{1}{2})$.}{$(\dfrac{3}{2};-1;\dfrac{1}{2})$.}{$(-\dfrac{5}{2};-1;\dfrac{1}{2})$.}

\cauTN[id=C2-B7-D06-NB-004,muc=NB,loigiai={Trung điểm $I$ của $AB$ có tọa độ $I\left(\dfrac{x_A+x_B}{2};\dfrac{y_A+y_B}{2};\dfrac{z_A+z_B}{2}\right)=(-1;-\dfrac{1}{2};1)$.}]{Trong không gian $Oxyz$, cho $A(-3;-3;2)$, $B(1;2;0)$. Tọa độ trung điểm $I$ của đoạn thẳng $AB$ là}
  {$(-1;-\dfrac{1}{2};1)$.}
  {$(-1;-\dfrac{5}{2};1)$.}{$(1;-\dfrac{1}{2};-1)$.}{$(-1;-\dfrac{1}{2};3)$.}

\cauTN[id=C2-B7-D06-NB-005,muc=NB,loigiai={Trung điểm $I$ của $AB$ có tọa độ $I\left(\dfrac{x_A+x_B}{2};\dfrac{y_A+y_B}{2};\dfrac{z_A+z_B}{2}\right)=(\dfrac{1}{2};-\dfrac{1}{2};0)$.}]{Trong không gian $Oxyz$, cho $A(0;-1;-1)$, $B(1;0;1)$. Tọa độ trung điểm $I$ của đoạn thẳng $AB$ là}
  {$(\dfrac{1}{2};-\dfrac{1}{2};0)$.}
  {$(\dfrac{1}{2};-\dfrac{3}{2};0)$.}{$(0;-\dfrac{1}{2};\dfrac{1}{2})$.}{$(1;1;2)$.}

\cauTN[id=C2-B7-D06-NB-006,muc=NB,loigiai={Trung điểm $I$ của $AB$ có tọa độ $I\left(\dfrac{x_A+x_B}{2};\dfrac{y_A+y_B}{2};\dfrac{z_A+z_B}{2}\right)=(-1;0;-\dfrac{3}{2})$.}]{Trong không gian $Oxyz$, cho $A(-3;0;-1)$, $B(1;0;-2)$. Tọa độ trung điểm $I$ của đoạn thẳng $AB$ là}
  {$(-1;0;-\dfrac{3}{2})$.}
  {$(-1;0;-\dfrac{5}{2})$.}{$(-2;0;-3)$.}{$(1;0;\dfrac{3}{2})$.}

\cauTN[id=C2-B7-D06-NB-007,muc=NB,loigiai={Trung điểm $I$ của $AB$ có tọa độ $I\left(\dfrac{x_A+x_B}{2};\dfrac{y_A+y_B}{2};\dfrac{z_A+z_B}{2}\right)=(-1;-1;-1)$.}]{Trong không gian $Oxyz$, cho $A(-3;-1;0)$, $B(1;-1;-2)$. Tọa độ trung điểm $I$ của đoạn thẳng $AB$ là}
  {$(-1;-1;-1)$.}
  {$(-3;-1;-1)$.}{$(1;1;1)$.}{$(-1;1;-1)$.}

\cauTN[id=C2-B7-D06-NB-008,muc=NB,loigiai={Trung điểm $I$ của $AB$ có tọa độ $I\left(\dfrac{x_A+x_B}{2};\dfrac{y_A+y_B}{2};\dfrac{z_A+z_B}{2}\right)=(2;\dfrac{5}{2};1)$.}]{Trong không gian $Oxyz$, cho $A(3;2;2)$, $B(1;3;0)$. Tọa độ trung điểm $I$ của đoạn thẳng $AB$ là}
  {$(2;\dfrac{5}{2};1)$.}
  {$(2;\dfrac{5}{2};0)$.}{$(\dfrac{5}{2};2;1)$.}{$(2;\dfrac{5}{2};3)$.}

\cauTN[id=C2-B7-D06-NB-009,muc=NB,loigiai={Trọng tâm $G$ của tam giác $ABC$ có $x_G=\dfrac{x_A+x_B+x_C}{3}$, tương tự cho $y_G$, $z_G$. Vậy $G(\dfrac{5}{3};2;-1)$.}]{Trong không gian $Oxyz$, cho $A(1;3;-1)$, $B(1;0;-2)$, $C(3;3;0)$. Tọa độ trọng tâm $G$ của tam giác $ABC$ là}
  {$(\dfrac{5}{3};2;-1)$.}
  {$(-\dfrac{1}{3};2;-1)$.}{$(\dfrac{8}{3};2;-1)$.}{$(\dfrac{2}{3};2;-1)$.}

\cauTN[id=C2-B7-D06-NB-010,muc=NB,loigiai={Trọng tâm $G$ của tam giác $ABC$ có $x_G=\dfrac{x_A+x_B+x_C}{3}$, tương tự cho $y_G$, $z_G$. Vậy $G(\dfrac{1}{3};-\dfrac{1}{3};-\dfrac{1}{3})$.}]{Trong không gian $Oxyz$, cho $A(1;3;1)$, $B(3;-3;-1)$, $C(-3;-1;-1)$. Tọa độ trọng tâm $G$ của tam giác $ABC$ là}
  {$(\dfrac{1}{3};-\dfrac{1}{3};-\dfrac{1}{3})$.}
  {$(\dfrac{1}{3};-\dfrac{1}{3};\dfrac{1}{3})$.}{$(\dfrac{1}{3};\dfrac{5}{3};-\dfrac{1}{3})$.}{$(\dfrac{1}{3};-\dfrac{7}{3};-\dfrac{1}{3})$.}

\cauTN[id=C2-B7-D06-NB-011,muc=NB,loigiai={Trọng tâm $G$ của tam giác $ABC$ có $x_G=\dfrac{x_A+x_B+x_C}{3}$, tương tự cho $y_G$, $z_G$. Vậy $G(0;\dfrac{1}{3};\dfrac{5}{3})$.}]{Trong không gian $Oxyz$, cho $A(1;2;3)$, $B(-2;1;1)$, $C(1;-2;1)$. Tọa độ trọng tâm $G$ của tam giác $ABC$ là}
  {$(0;\dfrac{1}{3};\dfrac{5}{3})$.}
  {$(\dfrac{5}{3};\dfrac{1}{3};0)$.}{$(0;1;5)$.}{$(0;\dfrac{1}{3};\dfrac{11}{3})$.}

\cauTN[id=C2-B7-D06-NB-012,muc=NB,loigiai={Trọng tâm $G$ của tam giác $ABC$ có $x_G=\dfrac{x_A+x_B+x_C}{3}$, tương tự cho $y_G$, $z_G$. Vậy $G(\dfrac{4}{3};-\dfrac{4}{3};\dfrac{5}{3})$.}]{Trong không gian $Oxyz$, cho $A(1;-2;2)$, $B(0;-2;3)$, $C(3;0;0)$. Tọa độ trọng tâm $G$ của tam giác $ABC$ là}
  {$(\dfrac{4}{3};-\dfrac{4}{3};\dfrac{5}{3})$.}
  {$(\dfrac{5}{3};-\dfrac{4}{3};\dfrac{4}{3})$.}{$(-\dfrac{4}{3};-\dfrac{4}{3};\dfrac{5}{3})$.}{$(\dfrac{4}{3};-\dfrac{7}{3};\dfrac{5}{3})$.}

\cauTN[id=C2-B7-D06-NB-013,muc=NB,loigiai={Trọng tâm $G$ của tam giác $ABC$ có $x_G=\dfrac{x_A+x_B+x_C}{3}$, tương tự cho $y_G$, $z_G$. Vậy $G(1;0;\dfrac{2}{3})$.}]{Trong không gian $Oxyz$, cho $A(-2;1;-2)$, $B(2;-2;1)$, $C(3;1;3)$. Tọa độ trọng tâm $G$ của tam giác $ABC$ là}
  {$(1;0;\dfrac{2}{3})$.}
  {$(1;2;\dfrac{2}{3})$.}{$(0;0;\dfrac{2}{3})$.}{$(-1;0;\dfrac{2}{3})$.}

\cauTN[id=C2-B7-D06-NB-014,muc=NB,loigiai={Trọng tâm $G$ của tam giác $ABC$ có $x_G=\dfrac{x_A+x_B+x_C}{3}$, tương tự cho $y_G$, $z_G$. Vậy $G(\dfrac{2}{3};\dfrac{1}{3};\dfrac{2}{3})$.}]{Trong không gian $Oxyz$, cho $A(0;-3;-1)$, $B(0;1;3)$, $C(2;3;0)$. Tọa độ trọng tâm $G$ của tam giác $ABC$ là}
  {$(\dfrac{2}{3};\dfrac{1}{3};\dfrac{2}{3})$.}
  {$(-\dfrac{2}{3};-\dfrac{1}{3};-\dfrac{2}{3})$.}{$(\dfrac{5}{3};\dfrac{1}{3};\dfrac{2}{3})$.}{$(2;1;2)$.}

\cauTN[id=C2-B7-D06-NB-015,muc=NB,loigiai={Trọng tâm $G$ của tam giác $ABC$ có $x_G=\dfrac{x_A+x_B+x_C}{3}$, tương tự cho $y_G$, $z_G$. Vậy $G(-\dfrac{1}{3};\dfrac{1}{3};1)$.}]{Trong không gian $Oxyz$, cho $A(-1;1;2)$, $B(-1;1;0)$, $C(1;-1;1)$. Tọa độ trọng tâm $G$ của tam giác $ABC$ là}
  {$(-\dfrac{1}{3};\dfrac{1}{3};1)$.}
  {$(-\dfrac{1}{3};\dfrac{1}{3};0)$.}{$(-\dfrac{1}{3};\dfrac{1}{3};2)$.}{$(1;\dfrac{1}{3};-\dfrac{1}{3})$.}

\cauTN[id=C2-B7-D06-NB-016,muc=NB,loigiai={Trọng tâm $G$ của tam giác $ABC$ có $x_G=\dfrac{x_A+x_B+x_C}{3}$, tương tự cho $y_G$, $z_G$. Vậy $G(\dfrac{1}{3};1;-\dfrac{2}{3})$.}]{Trong không gian $Oxyz$, cho $A(0;1;-1)$, $B(1;0;1)$, $C(0;2;-2)$. Tọa độ trọng tâm $G$ của tam giác $ABC$ là}
  {$(\dfrac{1}{3};1;-\dfrac{2}{3})$.}
  {$(\dfrac{1}{3};1;-\dfrac{5}{3})$.}{$(-\dfrac{2}{3};1;\dfrac{1}{3})$.}{$(\dfrac{1}{3};3;-\dfrac{2}{3})$.}

\cauTN[id=C2-B7-D06-NB-017,muc=NB,loigiai={Tứ giác $ABCD$ là hình bình hành khi $\overrightarrow{AB}=\overrightarrow{DC}$, suy ra $D=A+C-B=(3;2;-6)$.}]{Trong không gian $Oxyz$, cho $A(-2;1;-1)$, $B(-3;-2;2)$, $C(2;-1;-3)$. Tìm tọa độ điểm $D$ sao cho $ABCD$ là hình bình hành.}
  {$(3;2;-6)$.}
  {$(-7;0;4)$.}{$(3;2;-7)$.}{$(4;2;-6)$.}

\cauTN[id=C2-B7-D06-NB-018,muc=NB,loigiai={Tứ giác $ABCD$ là hình bình hành khi $\overrightarrow{AB}=\overrightarrow{DC}$, suy ra $D=A+C-B=(-2;-2;-1)$.}]{Trong không gian $Oxyz$, cho $A(-2;-3;-1)$, $B(-1;-2;-3)$, $C(-1;-1;-3)$. Tìm tọa độ điểm $D$ sao cho $ABCD$ là hình bình hành.}
  {$(-2;-2;-1)$.}
  {$(-2;-2;1)$.}{$(-1;-2;-1)$.}{$(-1;-2;-2)$.}

\cauTN[id=C2-B7-D06-NB-019,muc=NB,loigiai={Tứ giác $ABCD$ là hình bình hành khi $\overrightarrow{AB}=\overrightarrow{DC}$, suy ra $D=A+C-B=(-2;2;-2)$.}]{Trong không gian $Oxyz$, cho $A(-2;1;2)$, $B(3;-1;2)$, $C(3;0;-2)$. Tìm tọa độ điểm $D$ sao cho $ABCD$ là hình bình hành.}
  {$(-2;2;-2)$.}
  {$(-2;0;6)$.}{$(2;-2;-2)$.}{$(-2;1;-2)$.}

\cauTN[id=C2-B7-D06-NB-020,muc=NB,loigiai={Tứ giác $ABCD$ là hình bình hành khi $\overrightarrow{AB}=\overrightarrow{DC}$, suy ra $D=A+C-B=(6;2;-3)$.}]{Trong không gian $Oxyz$, cho $A(2;0;-2)$, $B(-3;-1;0)$, $C(1;1;-1)$. Tìm tọa độ điểm $D$ sao cho $ABCD$ là hình bình hành.}
  {$(6;2;-3)$.}
  {$(6;2;-5)$.}{$(6;-2;-3)$.}{$(-6;2;-3)$.}

\cauTN[id=C2-B7-D06-NB-021,muc=NB,loigiai={Tứ giác $ABCD$ là hình bình hành khi $\overrightarrow{AB}=\overrightarrow{DC}$, suy ra $D=A+C-B=(-4;7;5)$.}]{Trong không gian $Oxyz$, cho $A(-2;2;1)$, $B(0;-2;-3)$, $C(-2;3;1)$. Tìm tọa độ điểm $D$ sao cho $ABCD$ là hình bình hành.}
  {$(-4;7;5)$.}
  {$(0;-1;-3)$.}{$(-4;7;4)$.}{$(-4;7;3)$.}

\cauTN[id=C2-B7-D06-NB-022,muc=NB,loigiai={Tứ giác $ABCD$ là hình bình hành khi $\overrightarrow{AB}=\overrightarrow{DC}$, suy ra $D=A+C-B=(-2;2;-3)$.}]{Trong không gian $Oxyz$, cho $A(-1;-2;-1)$, $B(-2;-3;1)$, $C(-3;1;-1)$. Tìm tọa độ điểm $D$ sao cho $ABCD$ là hình bình hành.}
  {$(-2;2;-3)$.}
  {$(0;2;-3)$.}{$(-2;2;-2)$.}{$(-2;3;-3)$.}

\cauTN[id=C2-B7-D06-NB-023,muc=NB,loigiai={Từ giả thiết $\overrightarrow{AM}=\dfrac{2}{3}\overrightarrow{AB}$ nên $M=A+\dfrac{2}{3}(B-A)=(-\dfrac{4}{3};\dfrac{4}{3};-\dfrac{5}{3})$.}]{Trong không gian $Oxyz$, cho $A(0;0;-3)$, $B(-2;2;-1)$. Tọa độ điểm $M$ thuộc đoạn thẳng $AB$ sao cho $1\overrightarrow{AM}=2\overrightarrow{MB}$ là}
  {$(-\dfrac{4}{3};\dfrac{4}{3};-\dfrac{5}{3})$.}
  {$(-\dfrac{4}{3};-\dfrac{4}{3};-\dfrac{5}{3})$.}{$(-\dfrac{4}{3};\dfrac{4}{3};-\dfrac{8}{3})$.}{$(-\dfrac{10}{3};\dfrac{4}{3};-\dfrac{5}{3})$.}

\cauTN[id=C2-B7-D06-NB-024,muc=NB,loigiai={Từ giả thiết $\overrightarrow{AM}=\dfrac{3}{4}\overrightarrow{AB}$ nên $M=A+\dfrac{3}{4}(B-A)=(\dfrac{3}{2};-1;-3)$.}]{Trong không gian $Oxyz$, cho $A(0;2;-3)$, $B(2;-2;-3)$. Tọa độ điểm $M$ thuộc đoạn thẳng $AB$ sao cho $1\overrightarrow{AM}=3\overrightarrow{MB}$ là}
  {$(\dfrac{3}{2};-1;-3)$.}
  {$(\dfrac{3}{2};0;-3)$.}{$(-3;-1;\dfrac{3}{2})$.}{$(\dfrac{3}{2};-3;-3)$.}

\cauTN[id=C2-B7-D06-NB-025,muc=NB,loigiai={Từ giả thiết $\overrightarrow{AM}=-2\overrightarrow{AB}$ nên $M=A-2(B-A)=(-13;-9;-12)$.}]{Trong không gian $Oxyz$, cho $A(-3;-1;-2)$, $B(2;3;3)$. Tọa độ điểm $M$ thoả mãn $\overrightarrow{AM}=-2\overrightarrow{AB}$ là}
  {$(-13;-9;-12)$.}
  {$(-13;-11;-12)$.}{$(-9;-13;-12)$.}{$(-13;-7;-12)$.}

\cauTN[id=C2-B7-D06-NB-026,muc=NB,loigiai={Từ giả thiết $\overrightarrow{AM}=3\overrightarrow{AB}$ nên $M=A+3(B-A)=(-4;9;-4)$.}]{Trong không gian $Oxyz$, cho $A(-1;-3;-1)$, $B(-2;1;-2)$. Tọa độ điểm $M$ thoả mãn $\overrightarrow{AM}=3\overrightarrow{AB}$ là}
  {$(-4;9;-4)$.}
  {$(-4;9;4)$.}{$(1;-11;1)$.}{$(-4;9;-3)$.}

\cauTN[id=C2-B7-D06-NB-027,muc=NB,loigiai={Từ giả thiết $\overrightarrow{AM}=\dfrac{1}{4}\overrightarrow{AB}$ nên $M=A+\dfrac{1}{4}(B-A)=(\dfrac{5}{4};-1;-3)$.}]{Trong không gian $Oxyz$, cho $A(2;-2;-3)$, $B(-1;2;-3)$. Tọa độ điểm $M$ thuộc đoạn thẳng $AB$ sao cho $3\overrightarrow{AM}=1\overrightarrow{MB}$ là}
  {$(\dfrac{5}{4};-1;-3)$.}
  {$(\dfrac{13}{4};-1;-3)$.}{$(\dfrac{5}{4};0;-3)$.}{$(-3;-1;\dfrac{5}{4})$.}

\cauTN[id=C2-B7-D06-NB-028,muc=NB,loigiai={Từ giả thiết $\overrightarrow{AM}=-1\overrightarrow{AB}$ nên $M=A-1(B-A)=(-4;5;-6)$.}]{Trong không gian $Oxyz$, cho $A(-2;1;-2)$, $B(0;-3;2)$. Tọa độ điểm $M$ thoả mãn $\overrightarrow{AM}=-1\overrightarrow{AB}$ là}
  {$(-4;5;-6)$.}
  {$(-4;5;-5)$.}{$(-4;-5;-6)$.}{$(5;-4;-6)$.}

\cauTN[id=C2-B7-D06-NB-029,muc=NB,loigiai={Từ giả thiết $\overrightarrow{AM}=-2\overrightarrow{AB}$ nên $M=A-2(B-A)=(-3;11;-15)$.}]{Trong không gian $Oxyz$, cho $A(1;3;-3)$, $B(3;-1;3)$. Tọa độ điểm $M$ thoả mãn $\overrightarrow{AM}=-2\overrightarrow{AB}$ là}
  {$(-3;11;-15)$.}
  {$(-3;11;-14)$.}{$(-3;-11;-15)$.}{$(-5;11;-15)$.}

\cauTN[id=C2-B7-D06-NB-030,muc=NB,loigiai={Từ giả thiết $\overrightarrow{AM}=\dfrac{2}{3}\overrightarrow{AB}$ nên $M=A+\dfrac{2}{3}(B-A)=(-\dfrac{7}{3};-\dfrac{4}{3};-\dfrac{1}{3})$.}]{Trong không gian $Oxyz$, cho $A(-1;0;-1)$, $B(-3;-2;0)$. Tọa độ điểm $M$ thuộc đoạn thẳng $AB$ sao cho $1\overrightarrow{AM}=2\overrightarrow{MB}$ là}
  {$(-\dfrac{7}{3};-\dfrac{4}{3};-\dfrac{1}{3})$.}
  {$(-\dfrac{1}{3};-\dfrac{4}{3};-\dfrac{1}{3})$.}{$(-\dfrac{7}{3};-\dfrac{4}{3};\dfrac{5}{3})$.}{$(-4;-3;\dfrac{1}{2})$.}

\cauTN[id=C2-B7-D06-NB-031,muc=NB,loigiai={Vì $G$ là trọng tâm nên $C=3G-A-B=(1;1;1)$.}]{Cho tam giác $ABC$ có $A(-1;2;2)$, $B(1;3;1)$ và trọng tâm $G(\dfrac{1}{3};2;\dfrac{4}{3})$. Tọa độ điểm $C$ là}
  {$(1;1;1)$.}
  {$(1;1;-1)$.}{$(1;0;1)$.}{$(2;1;1)$.}

\cauTN[id=C2-B7-D06-NB-032,muc=NB,loigiai={Vì $G$ là trọng tâm nên $C=3G-A-B=(1;2;-1)$.}]{Cho tam giác $ABC$ có $A(2;1;-1)$, $B(2;1;-2)$ và trọng tâm $G(\dfrac{5}{3};\dfrac{4}{3};-\dfrac{4}{3})$. Tọa độ điểm $C$ là}
  {$(1;2;-1)$.}
  {$(-1;2;1)$.}{$(1;2;0)$.}{$(1;4;-1)$.}

\cauTN[id=C2-B7-D06-NB-033,muc=NB,loigiai={Vì $G$ là trọng tâm nên $C=3G-A-B=(-1;2;-1)$.}]{Cho tam giác $ABC$ có $A(3;0;3)$, $B(2;-3;3)$ và trọng tâm $G(\dfrac{4}{3};-\dfrac{1}{3};\dfrac{5}{3})$. Tọa độ điểm $C$ là}
  {$(-1;2;-1)$.}
  {$(0;2;-1)$.}{$(-1;0;-1)$.}{$(-1;2;0)$.}

\cauTN[id=C2-B7-D06-NB-034,muc=NB,loigiai={Vì $G$ là trọng tâm nên $C=3G-A-B=(-3;3;3)$.}]{Cho tam giác $ABC$ có $A(0;1;0)$, $B(-3;-3;-3)$ và trọng tâm $G(-2;\dfrac{1}{3};0)$. Tọa độ điểm $C$ là}
  {$(-3;3;3)$.}
  {$(-3;3;1)$.}{$(-3;3;-3)$.}{$(-3;5;3)$.}

\cauTN[id=C2-B7-D06-NB-035,muc=NB,loigiai={Vì $G$ là trọng tâm nên $C=3G-A-B=(0;3;-1)$.}]{Cho tam giác $ABC$ có $A(1;-2;1)$, $B(2;-3;-2)$ và trọng tâm $G(1;-\dfrac{2}{3};-\dfrac{2}{3})$. Tọa độ điểm $C$ là}
  {$(0;3;-1)$.}
  {$(0;-3;-1)$.}{$(2;3;-1)$.}{$(0;4;-1)$.}

\cauTN[id=C2-B7-D06-NB-036,muc=NB,loigiai={Vì $G$ là trọng tâm nên $C=3G-A-B=(-1;-2;2)$.}]{Cho tam giác $ABC$ có $A(3;0;-2)$, $B(-2;-2;-3)$ và trọng tâm $G(0;-\dfrac{4}{3};-1)$. Tọa độ điểm $C$ là}
  {$(-1;-2;2)$.}
  {$(-2;-2;2)$.}{$(-1;2;2)$.}{$(-1;-2;1)$.}

\cauDS[id=C2-B7-D06-NB-037,muc=NB,loigiai={Kiểm tra từng mệnh đề bằng công thức tọa độ. a) Đúng. b) Đúng. c) Sai. d) Đúng.}]{Trong không gian $Oxyz$, cho $A(-1;-3;-1)$, $B(2;-3;-3)$, $C(0;-3;1)$. Xét tính đúng sai của các mệnh đề sau:}{%
  \y{Trọng tâm tam giác $ABC$ có tọa độ $(\dfrac{1}{3};-3;-1)$.}{d}
  \y{Điểm $D$ để $ABCD$ là hình bình hành có tọa độ $(-3;-3;3)$.}{d}
  \y{Điểm $D$ để $ABCD$ là hình bình hành có tọa độ $(-3;-3;1)$.}{s}
  \y{Trung điểm của đoạn thẳng $AB$ có tọa độ $(\dfrac{1}{2};-3;-2)$.}{d}}

\cauDS[id=C2-B7-D06-NB-038,muc=NB,loigiai={Kiểm tra từng mệnh đề bằng công thức tọa độ. a) Sai. b) Sai. c) Đúng. d) Đúng.}]{Trong không gian $Oxyz$, cho $A(1;3;1)$, $B(-3;-3;-1)$, $C(1;-2;-2)$. Xét tính đúng sai của các mệnh đề sau:}{%
  \y{Tọa độ điểm $M$ thuộc đoạn thẳng $AB$ sao cho $1\overrightarrow{AM}=3\overrightarrow{MB}$ là $(-\dfrac{3}{2};-2;-\dfrac{1}{2})$.}{s}
  \y{Vectơ $\overrightarrow{AB}$ có tọa độ $(-4;-4;-2)$.}{s}
  \y{Điểm $D$ để $ABCD$ là hình bình hành có tọa độ $(5;4;0)$.}{d}
  \y{Trọng tâm tam giác $ABC$ có tọa độ $(-\dfrac{1}{3};-\dfrac{2}{3};-\dfrac{2}{3})$.}{d}}

\cauDS[id=C2-B7-D06-NB-039,muc=NB,loigiai={Kiểm tra từng mệnh đề bằng công thức tọa độ. a) Sai. b) Đúng. c) Đúng. d) Đúng.}]{Trong không gian $Oxyz$, cho $A(-3;2;2)$, $B(-1;2;-3)$, $C(-3;1;3)$. Xét tính đúng sai của các mệnh đề sau:}{%
  \y{Trọng tâm tam giác $ABC$ có tọa độ $(-\dfrac{7}{2};\dfrac{5}{2};1)$.}{s}
  \y{Tọa độ điểm $M$ thuộc đoạn thẳng $AB$ sao cho $3\overrightarrow{AM}=1\overrightarrow{MB}$ là $(-\dfrac{5}{2};2;\dfrac{3}{4})$.}{d}
  \y{Trung điểm của đoạn thẳng $AB$ có tọa độ $(-2;2;-\dfrac{1}{2})$.}{d}
  \y{Vectơ $\overrightarrow{AB}$ có tọa độ $(2;0;-5)$.}{d}}

\cauDS[id=C2-B7-D06-NB-040,muc=NB,loigiai={Kiểm tra từng mệnh đề bằng công thức tọa độ. a) Sai. b) Đúng. c) Đúng. d) Sai.}]{Trong không gian $Oxyz$, cho $A(2;3;-2)$, $B(0;-2;-2)$, $C(2;-3;-3)$. Xét tính đúng sai của các mệnh đề sau:}{%
  \y{Trung điểm của đoạn thẳng $AB$ có tọa độ $(1;-2;\dfrac{1}{2})$.}{s}
  \y{Trọng tâm tam giác $ABC$ có tọa độ $(\dfrac{4}{3};-\dfrac{2}{3};-\dfrac{7}{3})$.}{d}
  \y{Điểm $D$ để $ABCD$ là hình bình hành có tọa độ $(4;2;-3)$.}{d}
  \y{Tọa độ điểm $M$ thuộc đoạn thẳng $AB$ sao cho $2\overrightarrow{AM}=1\overrightarrow{MB}$ là $(\dfrac{4}{3};-\dfrac{4}{3};-2)$.}{s}}

\cauDS[id=C2-B7-D06-NB-041,muc=NB,loigiai={Kiểm tra từng mệnh đề bằng công thức tọa độ. a) Đúng. b) Sai. c) Đúng. d) Đúng.}]{Trong không gian $Oxyz$, cho $A(-3;2;-2)$, $B(3;1;-3)$, $C(0;2;0)$. Xét tính đúng sai của các mệnh đề sau:}{%
  \y{Vectơ $\overrightarrow{AB}$ có tọa độ $(6;-1;-1)$.}{d}
  \y{Tọa độ điểm $M$ thoả mãn $\overrightarrow{AM}=2\overrightarrow{AB}$ là $(9;0;-3)$.}{s}
  \y{Trọng tâm tam giác $ABC$ có tọa độ $(0;\dfrac{5}{3};-\dfrac{5}{3})$.}{d}
  \y{Điểm $D$ để $ABCD$ là hình bình hành có tọa độ $(-6;3;1)$.}{d}}

\cauDS[id=C2-B7-D06-NB-042,muc=NB,loigiai={Kiểm tra từng mệnh đề bằng công thức tọa độ. a) Sai. b) Sai. c) Đúng. d) Sai.}]{Trong không gian $Oxyz$, cho $A(0;2;-1)$, $B(0;-3;-2)$, $C(2;0;1)$. Xét tính đúng sai của các mệnh đề sau:}{%
  \y{Trọng tâm tam giác $ABC$ có tọa độ $(-\dfrac{1}{3};-\dfrac{1}{3};-\dfrac{2}{3})$.}{s}
  \y{Điểm $D$ để $ABCD$ là hình bình hành có tọa độ $(1;5;2)$.}{s}
  \y{Vectơ $\overrightarrow{AB}$ có tọa độ $(0;-5;-1)$.}{d}
  \y{Điểm $D$ để $ABCD$ là hình bình hành có tọa độ $(2;6;2)$.}{s}}

\cauDS[id=C2-B7-D06-NB-043,muc=NB,loigiai={Kiểm tra từng mệnh đề bằng công thức tọa độ. a) Đúng. b) Đúng. c) Sai. d) Sai.}]{Trong không gian $Oxyz$, cho $A(-3;1;3)$, $B(-2;3;2)$, $C(1;0;1)$. Xét tính đúng sai của các mệnh đề sau:}{%
  \y{Điểm $D$ để $ABCD$ là hình bình hành có tọa độ $(0;-2;2)$.}{d}
  \y{Vectơ $\overrightarrow{AB}$ có tọa độ $(1;2;-1)$.}{d}
  \y{Tọa độ điểm $M$ thoả mãn $\overrightarrow{AM}=3\overrightarrow{AB}$ là $(0;9;0)$.}{s}
  \y{Trung điểm của đoạn thẳng $AB$ có tọa độ $(-\dfrac{9}{2};2;\dfrac{5}{2})$.}{s}}

\cauDS[id=C2-B7-D06-NB-044,muc=NB,loigiai={Kiểm tra từng mệnh đề bằng công thức tọa độ. a) Đúng. b) Đúng. c) Sai. d) Sai.}]{Trong không gian $Oxyz$, cho $A(-1;-1;3)$, $B(-2;-2;0)$, $C(3;0;-3)$. Xét tính đúng sai của các mệnh đề sau:}{%
  \y{Điểm $D$ để $ABCD$ là hình bình hành có tọa độ $(4;1;0)$.}{d}
  \y{Trung điểm của đoạn thẳng $AB$ có tọa độ $(-\dfrac{3}{2};-\dfrac{3}{2};\dfrac{3}{2})$.}{d}
  \y{Vectơ $\overrightarrow{AB}$ có tọa độ $(1;-1;-3)$.}{s}
  \y{Trọng tâm tam giác $ABC$ có tọa độ $(0;-2;0)$.}{s}}

\cauDS[id=C2-B7-D06-NB-045,muc=NB,loigiai={Kiểm tra từng mệnh đề bằng công thức tọa độ. a) Sai. b) Đúng. c) Sai. d) Đúng.}]{Trong không gian $Oxyz$, cho $A(-2;3;-3)$, $B(1;-1;-3)$, $C(2;-1;1)$. Xét tính đúng sai của các mệnh đề sau:}{%
  \y{Điểm $D$ để $ABCD$ là hình bình hành có tọa độ $(1;-3;-1)$.}{s}
  \y{Tọa độ điểm $M$ thuộc đoạn thẳng $AB$ sao cho $2\overrightarrow{AM}=1\overrightarrow{MB}$ là $(-1;\dfrac{5}{3};-3)$.}{d}
  \y{Vectơ $\overrightarrow{AB}$ có tọa độ $(5;-4;0)$.}{s}
  \y{Điểm $D$ để $ABCD$ là hình bình hành có tọa độ $(-1;3;1)$.}{d}}

\cauDS[id=C2-B7-D06-NB-046,muc=NB,loigiai={Kiểm tra từng mệnh đề bằng công thức tọa độ. a) Đúng. b) Sai. c) Sai. d) Đúng.}]{Trong không gian $Oxyz$, cho $A(0;3;1)$, $B(-3;-2;3)$, $C(1;1;-3)$. Xét tính đúng sai của các mệnh đề sau:}{%
  \y{Vectơ $\overrightarrow{AB}$ có tọa độ $(-3;-5;2)$.}{d}
  \y{Trung điểm của đoạn thẳng $AB$ có tọa độ $(-3;1;4)$.}{s}
  \y{Trọng tâm tam giác $ABC$ có tọa độ $(-\dfrac{2}{3};\dfrac{2}{3};-\dfrac{2}{3})$.}{s}
  \y{Trọng tâm tam giác $ABC$ có tọa độ $(-\dfrac{2}{3};\dfrac{2}{3};\dfrac{1}{3})$.}{d}}

\cauDS[id=C2-B7-D06-NB-047,muc=NB,loigiai={Kiểm tra từng mệnh đề bằng công thức tọa độ. a) Đúng. b) Sai. c) Đúng. d) Sai.}]{Trong không gian $Oxyz$, cho $A(-3;2;0)$, $B(-3;2;-2)$, $C(2;-1;0)$. Xét tính đúng sai của các mệnh đề sau:}{%
  \y{Tọa độ điểm $M$ thuộc đoạn thẳng $AB$ sao cho $1\overrightarrow{AM}=2\overrightarrow{MB}$ là $(-3;2;-\dfrac{4}{3})$.}{d}
  \y{Tọa độ điểm $M$ thuộc đoạn thẳng $AB$ sao cho $2\overrightarrow{AM}=1\overrightarrow{MB}$ là $(-3;2;\dfrac{4}{3})$.}{s}
  \y{Trung điểm của đoạn thẳng $AB$ có tọa độ $(-3;2;-1)$.}{d}
  \y{Trọng tâm tam giác $ABC$ có tọa độ $(-\dfrac{4}{3};1;\dfrac{4}{3})$.}{s}}

\cauDS[id=C2-B7-D06-NB-048,muc=NB,loigiai={Kiểm tra từng mệnh đề bằng công thức tọa độ. a) Đúng. b) Sai. c) Đúng. d) Sai.}]{Trong không gian $Oxyz$, cho $A(1;3;-1)$, $B(-3;3;3)$, $C(1;0;3)$. Xét tính đúng sai của các mệnh đề sau:}{%
  \y{Điểm $D$ để $ABCD$ là hình bình hành có tọa độ $(5;0;-1)$.}{d}
  \y{Điểm $D$ để $ABCD$ là hình bình hành có tọa độ $(5;0;-2)$.}{s}
  \y{Trung điểm của đoạn thẳng $AB$ có tọa độ $(-1;3;1)$.}{d}
  \y{Trọng tâm tam giác $ABC$ có tọa độ $(-\dfrac{1}{3};2;-\dfrac{5}{3})$.}{s}}

\cauTLN[id=C2-B7-D06-NB-049,muc=NB,loigiai={Trung điểm $I$ của $AB$ có tọa độ $I\left(\dfrac{x_A+x_B}{2};\dfrac{y_A+y_B}{2};\dfrac{z_A+z_B}{2}\right)=(3;0;1)$.}]{Trong không gian $Oxyz$, cho $A(3;2;2)$, $B(3;-2;0)$. Gọi $I$ là trung điểm của đoạn thẳng $AB$. Tính $x_I+y_I+z_I$.}{4}

\cauTLN[id=C2-B7-D06-NB-050,muc=NB,loigiai={Trung điểm $I$ của $AB$ có tọa độ $I\left(\dfrac{x_A+x_B}{2};\dfrac{y_A+y_B}{2};\dfrac{z_A+z_B}{2}\right)=(-1;-2;-1)$.}]{Trong không gian $Oxyz$, cho $A(-3;-3;0)$, $B(1;-1;-2)$. Gọi $I$ là trung điểm của đoạn thẳng $AB$. Tính $x_I+y_I+z_I$.}{-4}

\cauTLN[id=C2-B7-D06-NB-051,muc=NB,loigiai={Trung điểm $I$ của $AB$ có tọa độ $I\left(\dfrac{x_A+x_B}{2};\dfrac{y_A+y_B}{2};\dfrac{z_A+z_B}{2}\right)=(1;1;0)$.}]{Trong không gian $Oxyz$, cho $A(1;-1;2)$, $B(1;3;-2)$. Gọi $I$ là trung điểm của đoạn thẳng $AB$. Tính $x_I+y_I+z_I$.}{2}

\cauTLN[id=C2-B7-D06-NB-052,muc=NB,loigiai={Trọng tâm $G$ của tam giác $ABC$ có $x_G=\dfrac{x_A+x_B+x_C}{3}$, tương tự cho $y_G$, $z_G$. Vậy $G(-2;1;-1)$.}]{Trong không gian $Oxyz$, cho $A(-3;0;-1)$, $B(-2;3;-1)$, $C(-1;0;-1)$. Gọi $G$ là trọng tâm tam giác $ABC$. Tính $x_G+y_G+z_G$.}{-2}

\cauTLN[id=C2-B7-D06-NB-053,muc=NB,loigiai={Trọng tâm $G$ của tam giác $ABC$ có $x_G=\dfrac{x_A+x_B+x_C}{3}$, tương tự cho $y_G$, $z_G$. Vậy $G(1;0;0)$.}]{Trong không gian $Oxyz$, cho $A(3;3;-1)$, $B(0;0;2)$, $C(0;-3;-1)$. Gọi $G$ là trọng tâm tam giác $ABC$. Tính $x_G+y_G+z_G$.}{1}

\cauTLN[id=C2-B7-D06-NB-054,muc=NB,loigiai={Trọng tâm $G$ của tam giác $ABC$ có $x_G=\dfrac{x_A+x_B+x_C}{3}$, tương tự cho $y_G$, $z_G$. Vậy $G(0;-3;-1)$.}]{Trong không gian $Oxyz$, cho $A(3;-3;0)$, $B(-3;-3;-2)$, $C(0;-3;-1)$. Gọi $G$ là trọng tâm tam giác $ABC$. Tính $x_G+y_G+z_G$.}{-4}

\cauTLN[id=C2-B7-D06-NB-055,muc=NB,loigiai={Tứ giác $ABCD$ là hình bình hành khi $\overrightarrow{AB}=\overrightarrow{DC}$, suy ra $D=A+C-B=(-3;8;5)$.}]{Trong không gian $Oxyz$, cho $A(-2;2;1)$, $B(2;-3;-3)$, $C(1;3;1)$. Tìm điểm $D$ sao cho $ABCD$ là hình bình hành. Tính $x_D+y_D+z_D$.}{10}

\cauTLN[id=C2-B7-D06-NB-056,muc=NB,loigiai={Tứ giác $ABCD$ là hình bình hành khi $\overrightarrow{AB}=\overrightarrow{DC}$, suy ra $D=A+C-B=(-5;-1;-3)$.}]{Trong không gian $Oxyz$, cho $A(-2;-1;-2)$, $B(3;1;-2)$, $C(0;1;-3)$. Tìm điểm $D$ sao cho $ABCD$ là hình bình hành. Tính $x_D+y_D+z_D$.}{-9}

\cauTLN[id=C2-B7-D06-NB-057,muc=NB,loigiai={Tứ giác $ABCD$ là hình bình hành khi $\overrightarrow{AB}=\overrightarrow{DC}$, suy ra $D=A+C-B=(-9;3;-4)$.}]{Trong không gian $Oxyz$, cho $A(-3;3;-3)$, $B(3;-2;1)$, $C(-3;-2;0)$. Tìm điểm $D$ sao cho $ABCD$ là hình bình hành. Tính $x_D+y_D+z_D$.}{-10}

\cauTLN[id=C2-B7-D06-NB-058,muc=NB,loigiai={Từ giả thiết $\overrightarrow{AM}=3\overrightarrow{AB}$ nên $M=A+3(B-A)=(-4;-15;9)$.}]{Trong không gian $Oxyz$, cho $A(-1;3;-3)$, $B(-2;-3;1)$. Tìm điểm $M$ thoả mãn $\overrightarrow{AM}=3\overrightarrow{AB}$. Tính $x_M+y_M+z_M$.}{-10}

\cauTLN[id=C2-B7-D06-NB-059,muc=NB,loigiai={Từ giả thiết $\overrightarrow{AM}=2\overrightarrow{AB}$ nên $M=A+2(B-A)=(6;3;-5)$.}]{Trong không gian $Oxyz$, cho $A(0;-1;3)$, $B(3;1;-1)$. Tìm điểm $M$ thoả mãn $\overrightarrow{AM}=2\overrightarrow{AB}$. Tính $x_M+y_M+z_M$.}{4}

\cauTLN[id=C2-B7-D06-NB-060,muc=NB,loigiai={Từ giả thiết $\overrightarrow{AM}=-2\overrightarrow{AB}$ nên $M=A-2(B-A)=(-11;3;6)$.}]{Trong không gian $Oxyz$, cho $A(-3;-1;2)$, $B(1;-3;0)$. Tìm điểm $M$ thoả mãn $\overrightarrow{AM}=-2\overrightarrow{AB}$. Tính $x_M+y_M+z_M$.}{-2}

\cauTN[id=C2-B7-D06-TH-001,muc=TH,loigiai={Vì $ABCD$ là hình bình hành nên $D=A+C-B=(0;-10;5)$. Tâm $I$ của hình bình hành là trung điểm của $AC$ (cũng là trung điểm $BD$): $I=(0;-3;1)$.}]{Trong không gian $Oxyz$, cho $A(-4;-5;1)$, $B(0;4;-3)$, $C(4;-1;1)$. Gọi $I$ là tâm của hình bình hành $ABCD$. Tọa độ điểm $I$ là}
  {$(0;-3;1)$.}
  {$(0;-4;1)$.}{$(0;-1;1)$.}{$(-3;0;1)$.}

\cauTN[id=C2-B7-D06-TH-002,muc=TH,loigiai={Vì $ABCD$ là hình bình hành nên $D=A+C-B=(-2;8;-1)$. Tâm $I$ của hình bình hành là trung điểm của $AC$ (cũng là trung điểm $BD$): $I=(-\dfrac{5}{2};\dfrac{3}{2};\dfrac{3}{2})$.}]{Trong không gian $Oxyz$, cho $A(-4;0;1)$, $B(-3;-5;4)$, $C(-1;3;2)$. Gọi $I$ là tâm của hình bình hành $ABCD$. Tọa độ điểm $I$ là}
  {$(-\dfrac{5}{2};\dfrac{3}{2};\dfrac{3}{2})$.}
  {$(-\dfrac{1}{2};\dfrac{3}{2};\dfrac{3}{2})$.}{$(-\dfrac{5}{2};\dfrac{1}{2};\dfrac{3}{2})$.}{$(-\dfrac{5}{2};\dfrac{7}{2};\dfrac{3}{2})$.}

\cauTN[id=C2-B7-D06-TH-003,muc=TH,loigiai={Vì $ABCD$ là hình bình hành nên $D=A+C-B=(3;-1;7)$. Tâm $I$ của hình bình hành là trung điểm của $AC$ (cũng là trung điểm $BD$): $I=(0;-\dfrac{3}{2};4)$.}]{Trong không gian $Oxyz$, cho $A(4;0;5)$, $B(-3;-2;1)$, $C(-4;-3;3)$. Gọi $I$ là tâm của hình bình hành $ABCD$. Tọa độ điểm $I$ là}
  {$(0;-\dfrac{3}{2};4)$.}
  {$(0;-\dfrac{3}{2};5)$.}{$(0;-\dfrac{5}{2};4)$.}{$(0;-\dfrac{3}{2};3)$.}

\cauTN[id=C2-B7-D06-TH-004,muc=TH,loigiai={Vì $ABCD$ là hình bình hành nên $D=A+C-B=(0;1;11)$. Tâm $I$ của hình bình hành là trung điểm của $AC$ (cũng là trung điểm $BD$): $I=(-2;-\dfrac{1}{2};\dfrac{7}{2})$.}]{Trong không gian $Oxyz$, cho $A(-5;3;2)$, $B(-4;-2;-4)$, $C(1;-4;5)$. Gọi $I$ là tâm của hình bình hành $ABCD$. Tọa độ điểm $I$ là}
  {$(-2;-\dfrac{1}{2};\dfrac{7}{2})$.}
  {$(-2;\dfrac{1}{2};\dfrac{7}{2})$.}{$(-2;-\dfrac{5}{2};\dfrac{7}{2})$.}{$(-2;-\dfrac{1}{2};\dfrac{11}{2})$.}

\cauTN[id=C2-B7-D06-TH-005,muc=TH,loigiai={Vì $ABCD$ là hình bình hành nên $D=A+C-B=(3;10;1)$. Tâm $I$ của hình bình hành là trung điểm của $AC$ (cũng là trung điểm $BD$): $I=(-1;\dfrac{5}{2};\dfrac{3}{2})$.}]{Trong không gian $Oxyz$, cho $A(-2;2;-1)$, $B(-5;-5;2)$, $C(0;3;4)$. Gọi $I$ là tâm của hình bình hành $ABCD$. Tọa độ điểm $I$ là}
  {$(-1;\dfrac{5}{2};\dfrac{3}{2})$.}
  {$(-1;\dfrac{5}{2};\dfrac{7}{2})$.}{$(-1;\dfrac{5}{2};\dfrac{1}{2})$.}{$(1;-\dfrac{5}{2};-\dfrac{3}{2})$.}

\cauTN[id=C2-B7-D06-TH-006,muc=TH,loigiai={Vì $ABCD$ là hình bình hành nên $D=A+C-B=(0;-6;4)$. Tâm $I$ của hình bình hành là trung điểm của $AC$ (cũng là trung điểm $BD$): $I=(2;-\dfrac{1}{2};1)$.}]{Trong không gian $Oxyz$, cho $A(4;3;4)$, $B(4;5;-2)$, $C(0;-4;-2)$. Gọi $I$ là tâm của hình bình hành $ABCD$. Tọa độ điểm $I$ là}
  {$(2;-\dfrac{1}{2};1)$.}
  {$(-2;\dfrac{1}{2};-1)$.}{$(-2;-\dfrac{1}{2};1)$.}{$(-\dfrac{1}{2};2;1)$.}

\cauTN[id=C2-B7-D06-TH-007,muc=TH,loigiai={Vì $ABCD$ là hình bình hành nên $D=A+C-B=(-2;2;5)$. Tâm $I$ của hình bình hành là trung điểm của $AC$ (cũng là trung điểm $BD$): $I=(\dfrac{1}{2};\dfrac{7}{2};1)$.}]{Trong không gian $Oxyz$, cho $A(-2;4;1)$, $B(3;5;-3)$, $C(3;3;1)$. Gọi $I$ là tâm của hình bình hành $ABCD$. Tọa độ điểm $I$ là}
  {$(\dfrac{1}{2};\dfrac{7}{2};1)$.}
  {$(\dfrac{7}{2};\dfrac{1}{2};1)$.}{$(\dfrac{1}{2};\dfrac{9}{2};1)$.}{$(\dfrac{1}{2};\dfrac{7}{2};0)$.}

\cauTN[id=C2-B7-D06-TH-008,muc=TH,loigiai={Vì $ABCD$ là hình bình hành nên $D=A+C-B=(4;8;12)$. Tâm $I$ của hình bình hành là trung điểm của $AC$ (cũng là trung điểm $BD$): $I=(\dfrac{1}{2};\dfrac{9}{2};4)$.}]{Trong không gian $Oxyz$, cho $A(2;5;4)$, $B(-3;1;-4)$, $C(-1;4;4)$. Gọi $I$ là tâm của hình bình hành $ABCD$. Tọa độ điểm $I$ là}
  {$(\dfrac{1}{2};\dfrac{9}{2};4)$.}
  {$(\dfrac{1}{2};\dfrac{11}{2};4)$.}{$(-\dfrac{1}{2};3;0)$.}{$(\dfrac{1}{2};4;\dfrac{9}{2})$.}

\cauTN[id=C2-B7-D06-TH-009,muc=TH,loigiai={Vì $ABCD$ là hình bình hành nên $D=A+C-B=(-10;5;-1)$. Tâm $I$ của hình bình hành là trung điểm của $AC$ (cũng là trung điểm $BD$): $I=(-\dfrac{7}{2};3;1)$.}]{Trong không gian $Oxyz$, cho $A(-5;3;2)$, $B(3;1;3)$, $C(-2;3;0)$. Gọi $I$ là tâm của hình bình hành $ABCD$. Tọa độ điểm $I$ là}
  {$(-\dfrac{7}{2};3;1)$.}
  {$(-\dfrac{7}{2};3;3)$.}{$(-\dfrac{7}{2};1;3)$.}{$(\dfrac{7}{2};-3;-1)$.}

\cauTN[id=C2-B7-D06-TH-010,muc=TH,loigiai={Vì $ABCD$ là hình bình hành nên $D=A+C-B=(6;1;1)$. Tâm $I$ của hình bình hành là trung điểm của $AC$ (cũng là trung điểm $BD$): $I=(\dfrac{5}{2};-\dfrac{3}{2};2)$.}]{Trong không gian $Oxyz$, cho $A(0;0;2)$, $B(-1;-4;3)$, $C(5;-3;2)$. Gọi $I$ là tâm của hình bình hành $ABCD$. Tọa độ điểm $I$ là}
  {$(\dfrac{5}{2};-\dfrac{3}{2};2)$.}
  {$(\dfrac{9}{2};-\dfrac{3}{2};2)$.}{$(-\dfrac{1}{2};-2;\dfrac{5}{2})$.}{$(\dfrac{5}{2};\dfrac{3}{2};2)$.}

\cauTN[id=C2-B7-D06-TH-011,muc=TH,loigiai={Vì $ABCD$ là hình bình hành nên $D=A+C-B=(11;-6;-6)$. Tâm $I$ của hình bình hành là trung điểm của $AC$ (cũng là trung điểm $BD$): $I=(3;-\dfrac{3}{2};-\dfrac{3}{2})$.}]{Trong không gian $Oxyz$, cho $A(4;1;-3)$, $B(-5;3;3)$, $C(2;-4;0)$. Gọi $I$ là tâm của hình bình hành $ABCD$. Tọa độ điểm $I$ là}
  {$(3;-\dfrac{3}{2};-\dfrac{3}{2})$.}
  {$(-\dfrac{1}{2};2;0)$.}{$(3;\dfrac{3}{2};-\dfrac{3}{2})$.}{$(3;-\dfrac{7}{2};-\dfrac{3}{2})$.}

\cauTN[id=C2-B7-D06-TH-012,muc=TH,loigiai={Vì $ABCD$ là hình bình hành nên $D=A+C-B=(4;5;-4)$. Tâm $I$ của hình bình hành là trung điểm của $AC$ (cũng là trung điểm $BD$): $I=(0;0;-\dfrac{3}{2})$.}]{Trong không gian $Oxyz$, cho $A(-1;-4;0)$, $B(-4;-5;1)$, $C(1;4;-3)$. Gọi $I$ là tâm của hình bình hành $ABCD$. Tọa độ điểm $I$ là}
  {$(0;0;-\dfrac{3}{2})$.}
  {$(-1;0;-\dfrac{3}{2})$.}{$(2;0;-\dfrac{3}{2})$.}{$(0;-2;-\dfrac{3}{2})$.}

\cauTN[id=C2-B7-D06-TH-013,muc=TH,loigiai={Trọng tâm tứ diện $ABCD$ có tọa độ là trung bình cộng tọa độ bốn đỉnh: $G=(\dfrac{7}{4};0;-\dfrac{9}{4})$.}]{Trong không gian $Oxyz$, cho $A(-2;0;0)$, $B(2;-2;-1)$, $C(4;4;-5)$, $D(3;-2;-3)$. Tọa độ trọng tâm $G$ của tứ diện $ABCD$ là}
  {$(\dfrac{7}{4};0;-\dfrac{9}{4})$.}
  {$(-\dfrac{9}{4};0;\dfrac{7}{4})$.}{$(\dfrac{7}{4};2;-\dfrac{9}{4})$.}{$(\dfrac{7}{4};0;\dfrac{9}{4})$.}

\cauTN[id=C2-B7-D06-TH-014,muc=TH,loigiai={Trọng tâm tứ diện $ABCD$ có tọa độ là trung bình cộng tọa độ bốn đỉnh: $G=(\dfrac{3}{4};-2;-\dfrac{7}{2})$.}]{Trong không gian $Oxyz$, cho $A(-2;-3;-1)$, $B(5;-3;-5)$, $C(-1;-3;-3)$, $D(1;1;-5)$. Tọa độ trọng tâm $G$ của tứ diện $ABCD$ là}
  {$(\dfrac{3}{4};-2;-\dfrac{7}{2})$.}
  {$(\dfrac{3}{4};-\dfrac{7}{2};-2)$.}{$(\dfrac{3}{4};-2;-\dfrac{9}{2})$.}{$(\dfrac{3}{4};-2;-\dfrac{5}{2})$.}

\cauTN[id=C2-B7-D06-TH-015,muc=TH,loigiai={Trọng tâm tứ diện $ABCD$ có tọa độ là trung bình cộng tọa độ bốn đỉnh: $G=(-2;2;\dfrac{7}{4})$.}]{Trong không gian $Oxyz$, cho $A(1;1;-2)$, $B(-1;5;5)$, $C(-3;1;2)$, $D(-5;1;2)$. Tọa độ trọng tâm $G$ của tứ diện $ABCD$ là}
  {$(-2;2;\dfrac{7}{4})$.}
  {$(-2;4;\dfrac{7}{4})$.}{$(2;-2;\dfrac{7}{4})$.}{$(-1;\dfrac{7}{3};\dfrac{5}{3})$.}

\cauTN[id=C2-B7-D06-TH-016,muc=TH,loigiai={Trọng tâm tứ diện $ABCD$ có tọa độ là trung bình cộng tọa độ bốn đỉnh: $G=(\dfrac{3}{4};\dfrac{9}{4};-\dfrac{1}{4})$.}]{Trong không gian $Oxyz$, cho $A(0;4;-5)$, $B(5;0;2)$, $C(2;3;-1)$, $D(-4;2;3)$. Tọa độ trọng tâm $G$ của tứ diện $ABCD$ là}
  {$(\dfrac{3}{4};\dfrac{9}{4};-\dfrac{1}{4})$.}
  {$(-\dfrac{1}{4};\dfrac{9}{4};\dfrac{3}{4})$.}{$(\dfrac{7}{3};\dfrac{7}{3};-\dfrac{4}{3})$.}{$(-\dfrac{3}{4};-\dfrac{9}{4};\dfrac{1}{4})$.}

\cauTN[id=C2-B7-D06-TH-017,muc=TH,loigiai={Trọng tâm tứ diện $ABCD$ có tọa độ là trung bình cộng tọa độ bốn đỉnh: $G=(-\dfrac{9}{4};\dfrac{7}{4};3)$.}]{Trong không gian $Oxyz$, cho $A(-2;5;1)$, $B(-5;1;4)$, $C(0;3;2)$, $D(-2;-2;5)$. Tọa độ trọng tâm $G$ của tứ diện $ABCD$ là}
  {$(-\dfrac{9}{4};\dfrac{7}{4};3)$.}
  {$(-\dfrac{9}{4};\dfrac{7}{4};5)$.}{$(-\dfrac{9}{4};\dfrac{11}{4};3)$.}{$(-\dfrac{9}{4};-\dfrac{1}{4};3)$.}

\cauTN[id=C2-B7-D06-TH-018,muc=TH,loigiai={Trọng tâm tứ diện $ABCD$ có tọa độ là trung bình cộng tọa độ bốn đỉnh: $G=(\dfrac{1}{4};-\dfrac{3}{2};-2)$.}]{Trong không gian $Oxyz$, cho $A(-2;4;-5)$, $B(-2;-5;-1)$, $C(3;-1;-2)$, $D(2;-4;0)$. Tọa độ trọng tâm $G$ của tứ diện $ABCD$ là}
  {$(\dfrac{1}{4};-\dfrac{3}{2};-2)$.}
  {$(-\dfrac{1}{4};\dfrac{3}{2};2)$.}{$(\dfrac{1}{4};-\dfrac{1}{2};-2)$.}{$(\dfrac{5}{4};-\dfrac{3}{2};-2)$.}

\cauTN[id=C2-B7-D06-TH-019,muc=TH,loigiai={Trọng tâm tứ diện $ABCD$ có tọa độ là trung bình cộng tọa độ bốn đỉnh: $G=(0;\dfrac{1}{2};-\dfrac{3}{4})$.}]{Trong không gian $Oxyz$, cho $A(4;3;-3)$, $B(-2;-1;2)$, $C(-1;-2;2)$, $D(-1;2;-4)$. Tọa độ trọng tâm $G$ của tứ diện $ABCD$ là}
  {$(0;\dfrac{1}{2};-\dfrac{3}{4})$.}
  {$(0;\dfrac{1}{2};\dfrac{1}{4})$.}{$(0;-\dfrac{3}{2};-\dfrac{3}{4})$.}{$(\dfrac{1}{3};0;\dfrac{1}{3})$.}

\cauTN[id=C2-B7-D06-TH-020,muc=TH,loigiai={Trọng tâm tứ diện $ABCD$ có tọa độ là trung bình cộng tọa độ bốn đỉnh: $G=(\dfrac{1}{4};-\dfrac{5}{4};\dfrac{9}{4})$.}]{Trong không gian $Oxyz$, cho $A(4;-1;-2)$, $B(0;3;5)$, $C(-2;-5;2)$, $D(-1;-2;4)$. Tọa độ trọng tâm $G$ của tứ diện $ABCD$ là}
  {$(\dfrac{1}{4};-\dfrac{5}{4};\dfrac{9}{4})$.}
  {$(\dfrac{1}{4};-\dfrac{13}{4};\dfrac{9}{4})$.}{$(\dfrac{1}{4};-\dfrac{1}{4};\dfrac{9}{4})$.}{$(\dfrac{9}{4};-\dfrac{5}{4};\dfrac{9}{4})$.}

\cauTN[id=C2-B7-D06-TH-021,muc=TH,loigiai={Trọng tâm tứ diện $ABCD$ có tọa độ là trung bình cộng tọa độ bốn đỉnh: $G=(\dfrac{5}{2};-1;\dfrac{3}{2})$.}]{Trong không gian $Oxyz$, cho $A(5;1;2)$, $B(5;-4;4)$, $C(2;-1;4)$, $D(-2;0;-4)$. Tọa độ trọng tâm $G$ của tứ diện $ABCD$ là}
  {$(\dfrac{5}{2};-1;\dfrac{3}{2})$.}
  {$(-\dfrac{5}{2};-1;\dfrac{3}{2})$.}{$(\dfrac{5}{2};1;\dfrac{3}{2})$.}{$(\dfrac{1}{2};-1;\dfrac{3}{2})$.}

\cauTN[id=C2-B7-D06-TH-022,muc=TH,loigiai={Trọng tâm tứ diện $ABCD$ có tọa độ là trung bình cộng tọa độ bốn đỉnh: $G=(-4;-1;0)$.}]{Trong không gian $Oxyz$, cho $A(-5;-3;-3)$, $B(-5;-1;-1)$, $C(-1;-1;5)$, $D(-5;1;-1)$. Tọa độ trọng tâm $G$ của tứ diện $ABCD$ là}
  {$(-4;-1;0)$.}
  {$(-2;-1;0)$.}{$(-4;-1;2)$.}{$(-3;-1;0)$.}

\cauTN[id=C2-B7-D06-TH-023,muc=TH,loigiai={Trọng tâm tứ diện $ABCD$ có tọa độ là trung bình cộng tọa độ bốn đỉnh: $G=(2;-1;\dfrac{7}{2})$.}]{Trong không gian $Oxyz$, cho $A(3;-2;4)$, $B(-2;-5;2)$, $C(4;0;5)$, $D(3;3;3)$. Tọa độ trọng tâm $G$ của tứ diện $ABCD$ là}
  {$(2;-1;\dfrac{7}{2})$.}
  {$(2;-2;\dfrac{7}{2})$.}{$(2;\dfrac{7}{2};-1)$.}{$(2;-1;\dfrac{9}{2})$.}

\cauTN[id=C2-B7-D06-TH-024,muc=TH,loigiai={Trọng tâm tứ diện $ABCD$ có tọa độ là trung bình cộng tọa độ bốn đỉnh: $G=(-\dfrac{11}{4};-\dfrac{5}{2};-\dfrac{5}{4})$.}]{Trong không gian $Oxyz$, cho $A(-2;-5;-1)$, $B(-3;5;-1)$, $C(-2;-5;-4)$, $D(-4;-5;1)$. Tọa độ trọng tâm $G$ của tứ diện $ABCD$ là}
  {$(-\dfrac{11}{4};-\dfrac{5}{2};-\dfrac{5}{4})$.}
  {$(-\dfrac{11}{4};-\dfrac{5}{2};-\dfrac{1}{4})$.}{$(-\dfrac{11}{4};\dfrac{5}{2};-\dfrac{5}{4})$.}{$(-\dfrac{5}{2};-\dfrac{11}{4};-\dfrac{5}{4})$.}

\cauTN[id=C2-B7-D06-TH-025,muc=TH,loigiai={Gọi $M$, $N$ lần lượt là trung điểm $AB$, $AC$: $M=(0;1;0)$, $N=(0;\dfrac{1}{2};-\dfrac{1}{2})$ nên $\overrightarrow{MN}=(0;-\dfrac{1}{2};-\dfrac{1}{2})$.}]{Trong không gian $Oxyz$, cho $A(4;-2;-3)$, $B(-4;4;3)$, $C(-4;3;2)$. Gọi $M$, $N$ lần lượt là trung điểm của $AB$ và $AC$. Tọa độ vectơ $\overrightarrow{MN}$ là}
  {$(0;-\dfrac{1}{2};-\dfrac{1}{2})$.}
  {$(0;-\dfrac{1}{2};-\dfrac{3}{2})$.}{$(0;-\dfrac{1}{2};\dfrac{1}{2})$.}{$(0;1;1)$.}

\cauTN[id=C2-B7-D06-TH-026,muc=TH,loigiai={Gọi $M$, $N$ lần lượt là trung điểm $AB$, $AC$: $M=(-\dfrac{1}{2};-3;\dfrac{7}{2})$, $N=(2;-2;2)$ nên $\overrightarrow{MN}=(\dfrac{5}{2};1;-\dfrac{3}{2})$.}]{Trong không gian $Oxyz$, cho $A(-1;-4;5)$, $B(0;-2;2)$, $C(5;0;-1)$. Gọi $M$, $N$ lần lượt là trung điểm của $AB$ và $AC$. Tọa độ vectơ $\overrightarrow{MN}$ là}
  {$(\dfrac{5}{2};1;-\dfrac{3}{2})$.}
  {$(\dfrac{1}{2};1;-\dfrac{3}{2})$.}{$(\dfrac{7}{2};1;-\dfrac{3}{2})$.}{$(\dfrac{5}{2};1;\dfrac{3}{2})$.}

\cauTN[id=C2-B7-D06-TH-027,muc=TH,loigiai={Gọi $M$, $N$ lần lượt là trung điểm $AB$, $AC$: $M=(-1;-1;\dfrac{3}{2})$, $N=(-2;-1;1)$ nên $\overrightarrow{MN}=(-1;0;-\dfrac{1}{2})$.}]{Trong không gian $Oxyz$, cho $A(1;1;4)$, $B(-3;-3;-1)$, $C(-5;-3;-2)$. Gọi $M$, $N$ lần lượt là trung điểm của $AB$ và $AC$. Tọa độ vectơ $\overrightarrow{MN}$ là}
  {$(-1;0;-\dfrac{1}{2})$.}
  {$(-2;0;-\dfrac{1}{2})$.}{$(-1;2;-\dfrac{1}{2})$.}{$(-1;-\dfrac{1}{2};0)$.}

\cauTN[id=C2-B7-D06-TH-028,muc=TH,loigiai={Gọi $M$, $N$ lần lượt là trung điểm $AB$, $AC$: $M=(\dfrac{1}{2};\dfrac{3}{2};-\dfrac{5}{2})$, $N=(0;0;0)$ nên $\overrightarrow{MN}=(-\dfrac{1}{2};-\dfrac{3}{2};\dfrac{5}{2})$.}]{Trong không gian $Oxyz$, cho $A(0;-1;-4)$, $B(1;4;-1)$, $C(0;1;4)$. Gọi $M$, $N$ lần lượt là trung điểm của $AB$ và $AC$. Tọa độ vectơ $\overrightarrow{MN}$ là}
  {$(-\dfrac{1}{2};-\dfrac{3}{2};\dfrac{5}{2})$.}
  {$(-\dfrac{1}{2};-\dfrac{7}{2};\dfrac{5}{2})$.}{$(-\dfrac{1}{2};-\dfrac{3}{2};\dfrac{7}{2})$.}{$(-\dfrac{5}{2};-\dfrac{3}{2};\dfrac{5}{2})$.}

\cauTN[id=C2-B7-D06-TH-029,muc=TH,loigiai={Gọi $M$, $N$ lần lượt là trung điểm $AB$, $AC$: $M=(\dfrac{1}{2};2;1)$, $N=(\dfrac{3}{2};-3;2)$ nên $\overrightarrow{MN}=(1;-5;1)$.}]{Trong không gian $Oxyz$, cho $A(-2;-1;0)$, $B(3;5;2)$, $C(5;-5;4)$. Gọi $M$, $N$ lần lượt là trung điểm của $AB$ và $AC$. Tọa độ vectơ $\overrightarrow{MN}$ là}
  {$(1;-5;1)$.}
  {$(1;-5;-1)$.}{$(1;-7;1)$.}{$(1;-3;1)$.}

\cauTN[id=C2-B7-D06-TH-030,muc=TH,loigiai={Gọi $M$, $N$ lần lượt là trung điểm $AB$, $AC$: $M=(-4;1;0)$, $N=(-3;-\dfrac{1}{2};\dfrac{3}{2})$ nên $\overrightarrow{MN}=(1;-\dfrac{3}{2};\dfrac{3}{2})$.}]{Trong không gian $Oxyz$, cho $A(-5;4;-1)$, $B(-3;-2;1)$, $C(-1;-5;4)$. Gọi $M$, $N$ lần lượt là trung điểm của $AB$ và $AC$. Tọa độ vectơ $\overrightarrow{MN}$ là}
  {$(1;-\dfrac{3}{2};\dfrac{3}{2})$.}
  {$(1;\dfrac{1}{2};\dfrac{3}{2})$.}{$(1;-\dfrac{1}{2};\dfrac{3}{2})$.}{$(2;-\dfrac{3}{2};\dfrac{3}{2})$.}

\cauTN[id=C2-B7-D06-TH-031,muc=TH,loigiai={Gọi $M$, $N$ lần lượt là trung điểm $AB$, $AC$: $M=(-\dfrac{1}{2};-\dfrac{1}{2};\dfrac{3}{2})$, $N=(1;\dfrac{1}{2};3)$ nên $\overrightarrow{MN}=(\dfrac{3}{2};1;\dfrac{3}{2})$.}]{Trong không gian $Oxyz$, cho $A(1;2;3)$, $B(-2;-3;0)$, $C(1;-1;3)$. Gọi $M$, $N$ lần lượt là trung điểm của $AB$ và $AC$. Tọa độ vectơ $\overrightarrow{MN}$ là}
  {$(\dfrac{3}{2};1;\dfrac{3}{2})$.}
  {$(\dfrac{1}{2};1;\dfrac{3}{2})$.}{$(\dfrac{3}{2};-1;\dfrac{3}{2})$.}{$(\dfrac{5}{2};1;\dfrac{3}{2})$.}

\cauTN[id=C2-B7-D06-TH-032,muc=TH,loigiai={Gọi $M$, $N$ lần lượt là trung điểm $AB$, $AC$: $M=(\dfrac{3}{2};1;\dfrac{1}{2})$, $N=(-\dfrac{1}{2};-1;\dfrac{3}{2})$ nên $\overrightarrow{MN}=(-2;-2;1)$.}]{Trong không gian $Oxyz$, cho $A(4;3;0)$, $B(-1;-1;1)$, $C(-5;-5;3)$. Gọi $M$, $N$ lần lượt là trung điểm của $AB$ và $AC$. Tọa độ vectơ $\overrightarrow{MN}$ là}
  {$(-2;-2;1)$.}
  {$(-2;-2;-1)$.}{$(-2;2;1)$.}{$(-2;-4;1)$.}

\cauTN[id=C2-B7-D06-TH-033,muc=TH,loigiai={Gọi $M$, $N$ lần lượt là trung điểm $AB$, $AC$: $M=(-2;\dfrac{1}{2};0)$, $N=(\dfrac{3}{2};-3;-\dfrac{5}{2})$ nên $\overrightarrow{MN}=(\dfrac{7}{2};-\dfrac{7}{2};-\dfrac{5}{2})$.}]{Trong không gian $Oxyz$, cho $A(-1;-1;0)$, $B(-3;2;0)$, $C(4;-5;-5)$. Gọi $M$, $N$ lần lượt là trung điểm của $AB$ và $AC$. Tọa độ vectơ $\overrightarrow{MN}$ là}
  {$(\dfrac{7}{2};-\dfrac{7}{2};-\dfrac{5}{2})$.}
  {$(\dfrac{9}{2};-\dfrac{7}{2};-\dfrac{5}{2})$.}{$(-\dfrac{7}{2};-\dfrac{7}{2};-\dfrac{5}{2})$.}{$(\dfrac{11}{2};-\dfrac{7}{2};-\dfrac{5}{2})$.}

\cauTN[id=C2-B7-D06-TH-034,muc=TH,loigiai={Gọi $M$, $N$ lần lượt là trung điểm $AB$, $AC$: $M=(-3;3;-\dfrac{1}{2})$, $N=(-\dfrac{1}{2};0;-2)$ nên $\overrightarrow{MN}=(\dfrac{5}{2};-3;-\dfrac{3}{2})$.}]{Trong không gian $Oxyz$, cho $A(-1;2;-1)$, $B(-5;4;0)$, $C(0;-2;-3)$. Gọi $M$, $N$ lần lượt là trung điểm của $AB$ và $AC$. Tọa độ vectơ $\overrightarrow{MN}$ là}
  {$(\dfrac{5}{2};-3;-\dfrac{3}{2})$.}
  {$(\dfrac{5}{2};-3;-\dfrac{7}{2})$.}{$(-5;6;3)$.}{$(-3;\dfrac{5}{2};-\dfrac{3}{2})$.}

\cauTN[id=C2-B7-D06-TH-035,muc=TH,loigiai={Gọi $M$, $N$ lần lượt là trung điểm $AB$, $AC$: $M=(\dfrac{1}{2};-\dfrac{3}{2};\dfrac{3}{2})$, $N=(\dfrac{9}{2};-\dfrac{1}{2};\dfrac{3}{2})$ nên $\overrightarrow{MN}=(4;1;0)$.}]{Trong không gian $Oxyz$, cho $A(4;1;2)$, $B(-3;-4;1)$, $C(5;-2;1)$. Gọi $M$, $N$ lần lượt là trung điểm của $AB$ và $AC$. Tọa độ vectơ $\overrightarrow{MN}$ là}
  {$(4;1;0)$.}
  {$(5;1;0)$.}{$(4;-1;0)$.}{$(2;1;0)$.}

\cauTN[id=C2-B7-D06-TH-036,muc=TH,loigiai={Gọi $M$, $N$ lần lượt là trung điểm $AB$, $AC$: $M=(-1;\dfrac{7}{2};1)$, $N=(2;\dfrac{7}{2};3)$ nên $\overrightarrow{MN}=(3;0;2)$.}]{Trong không gian $Oxyz$, cho $A(3;2;2)$, $B(-5;5;0)$, $C(1;5;4)$. Gọi $M$, $N$ lần lượt là trung điểm của $AB$ và $AC$. Tọa độ vectơ $\overrightarrow{MN}$ là}
  {$(3;0;2)$.}
  {$(3;-2;2)$.}{$(3;0;3)$.}{$(3;2;2)$.}

\cauTN[id=C2-B7-D06-TH-037,muc=TH,loigiai={Từ $\overrightarrow{MA}=2\overrightarrow{MB}$ suy ra $A-M=2(B-M)$, do đó $M=\dfrac{A-2B}{1-2}=(15;-1;2)$.}]{Trong không gian $Oxyz$, cho $A(-5;1;-2)$, $B(5;0;0)$. Tìm tọa độ điểm $M$ thoả mãn $\overrightarrow{MA}=2\overrightarrow{MB}$.}
  {$(15;-1;2)$.}
  {$(17;-1;2)$.}{$(-1;15;2)$.}{$(15;-1;-2)$.}

\cauTN[id=C2-B7-D06-TH-038,muc=TH,loigiai={Từ $\overrightarrow{MA}=-1\overrightarrow{MB}$ suy ra $A-M=-1(B-M)$, do đó $M=\dfrac{A+1B}{1+1}=(\dfrac{1}{2};0;0)$.}]{Trong không gian $Oxyz$, cho $A(2;-5;-2)$, $B(-1;5;2)$. Tìm tọa độ điểm $M$ thoả mãn $\overrightarrow{MA}=-1\overrightarrow{MB}$.}
  {$(\dfrac{1}{2};0;0)$.}
  {$(-\dfrac{1}{2};0;0)$.}{$(\dfrac{1}{2};0;2)$.}{$(\dfrac{1}{2};-2;0)$.}

\cauTN[id=C2-B7-D06-TH-039,muc=TH,loigiai={Từ $\overrightarrow{MA}=-\dfrac{1}{2}\overrightarrow{MB}$ suy ra $A-M=-\dfrac{1}{2}(B-M)$, do đó $M=\dfrac{A+\dfrac{1}{2}B}{1+\dfrac{1}{2}}=(-\dfrac{5}{3};-\dfrac{11}{3};\dfrac{5}{3})$.}]{Trong không gian $Oxyz$, cho $A(-1;-5;3)$, $B(-3;-1;-1)$. Tìm tọa độ điểm $M$ thoả mãn $\overrightarrow{MA}=-\dfrac{1}{2}\overrightarrow{MB}$.}
  {$(-\dfrac{5}{3};-\dfrac{11}{3};\dfrac{5}{3})$.}
  {$(-\dfrac{5}{3};-\dfrac{11}{3};\dfrac{8}{3})$.}{$(-\dfrac{5}{3};-\dfrac{14}{3};\dfrac{5}{3})$.}{$(-\dfrac{8}{3};-\dfrac{11}{3};\dfrac{5}{3})$.}

\cauTN[id=C2-B7-D06-TH-040,muc=TH,loigiai={Từ $\overrightarrow{MA}=2\overrightarrow{MB}$ suy ra $A-M=2(B-M)$, do đó $M=\dfrac{A-2B}{1-2}=(1;2;-3)$.}]{Trong không gian $Oxyz$, cho $A(-3;-2;-5)$, $B(-1;0;-4)$. Tìm tọa độ điểm $M$ thoả mãn $\overrightarrow{MA}=2\overrightarrow{MB}$.}
  {$(1;2;-3)$.}
  {$(1;2;-1)$.}{$(-1;2;-3)$.}{$(-3;2;1)$.}

\cauTN[id=C2-B7-D06-TH-041,muc=TH,loigiai={Từ $\overrightarrow{MA}=-1\overrightarrow{MB}$ suy ra $A-M=-1(B-M)$, do đó $M=\dfrac{A+1B}{1+1}=(\dfrac{7}{2};-4;5)$.}]{Trong không gian $Oxyz$, cho $A(2;-3;5)$, $B(5;-5;5)$. Tìm tọa độ điểm $M$ thoả mãn $\overrightarrow{MA}=-1\overrightarrow{MB}$.}
  {$(\dfrac{7}{2};-4;5)$.}
  {$(-\dfrac{7}{2};4;-5)$.}{$(-\dfrac{7}{2};-4;5)$.}{$(-4;\dfrac{7}{2};5)$.}

\cauTN[id=C2-B7-D06-TH-042,muc=TH,loigiai={Từ $\overrightarrow{MA}=3\overrightarrow{MB}$ suy ra $A-M=3(B-M)$, do đó $M=\dfrac{A-3B}{1-3}=(-\dfrac{11}{2};\dfrac{1}{2};\dfrac{7}{2})$.}]{Trong không gian $Oxyz$, cho $A(5;-1;-1)$, $B(-2;0;2)$. Tìm tọa độ điểm $M$ thoả mãn $\overrightarrow{MA}=3\overrightarrow{MB}$.}
  {$(-\dfrac{11}{2};\dfrac{1}{2};\dfrac{7}{2})$.}
  {$(\dfrac{11}{2};\dfrac{1}{2};\dfrac{7}{2})$.}{$(-\dfrac{11}{2};\dfrac{7}{2};\dfrac{1}{2})$.}{$(-\dfrac{11}{2};\dfrac{1}{2};\dfrac{9}{2})$.}

\cauTN[id=C2-B7-D06-TH-043,muc=TH,loigiai={Từ $\overrightarrow{MA}=\dfrac{1}{2}\overrightarrow{MB}$ suy ra $A-M=\dfrac{1}{2}(B-M)$, do đó $M=\dfrac{A-\dfrac{1}{2}B}{1-\dfrac{1}{2}}=(5;6;3)$.}]{Trong không gian $Oxyz$, cho $A(1;4;2)$, $B(-3;2;1)$. Tìm tọa độ điểm $M$ thoả mãn $\overrightarrow{MA}=\dfrac{1}{2}\overrightarrow{MB}$.}
  {$(5;6;3)$.}
  {$(5;6;2)$.}{$(5;6;5)$.}{$(5;6;-3)$.}

\cauTN[id=C2-B7-D06-TH-044,muc=TH,loigiai={Từ $\overrightarrow{MA}=3\overrightarrow{MB}$ suy ra $A-M=3(B-M)$, do đó $M=\dfrac{A-3B}{1-3}=(\dfrac{7}{2};\dfrac{9}{2};-1)$.}]{Trong không gian $Oxyz$, cho $A(-4;0;2)$, $B(1;3;0)$. Tìm tọa độ điểm $M$ thoả mãn $\overrightarrow{MA}=3\overrightarrow{MB}$.}
  {$(\dfrac{7}{2};\dfrac{9}{2};-1)$.}
  {$(\dfrac{7}{2};\dfrac{11}{2};-1)$.}{$(-1;\dfrac{9}{2};\dfrac{7}{2})$.}{$(-\dfrac{7}{2};-\dfrac{9}{2};1)$.}

\cauTN[id=C2-B7-D06-TH-045,muc=TH,loigiai={Từ $\overrightarrow{MA}=-2\overrightarrow{MB}$ suy ra $A-M=-2(B-M)$, do đó $M=\dfrac{A+2B}{1+2}=(\dfrac{1}{3};-\dfrac{4}{3};\dfrac{4}{3})$.}]{Trong không gian $Oxyz$, cho $A(-5;-2;4)$, $B(3;-1;0)$. Tìm tọa độ điểm $M$ thoả mãn $\overrightarrow{MA}=-2\overrightarrow{MB}$.}
  {$(\dfrac{1}{3};-\dfrac{4}{3};\dfrac{4}{3})$.}
  {$(\dfrac{1}{3};\dfrac{2}{3};\dfrac{4}{3})$.}{$(\dfrac{1}{3};-\dfrac{4}{3};\dfrac{7}{3})$.}{$(\dfrac{4}{3};-\dfrac{4}{3};\dfrac{4}{3})$.}

\cauTN[id=C2-B7-D06-TH-046,muc=TH,loigiai={Từ $\overrightarrow{MA}=\dfrac{1}{2}\overrightarrow{MB}$ suy ra $A-M=\dfrac{1}{2}(B-M)$, do đó $M=\dfrac{A-\dfrac{1}{2}B}{1-\dfrac{1}{2}}=(-3;-13;6)$.}]{Trong không gian $Oxyz$, cho $A(-3;-5;2)$, $B(-3;3;-2)$. Tìm tọa độ điểm $M$ thoả mãn $\overrightarrow{MA}=\dfrac{1}{2}\overrightarrow{MB}$.}
  {$(-3;-13;6)$.}
  {$(-3;-13;8)$.}{$(-3;-13;5)$.}{$(-3;-14;6)$.}

\cauTN[id=C2-B7-D06-TH-047,muc=TH,loigiai={Từ $\overrightarrow{MA}=-\dfrac{1}{2}\overrightarrow{MB}$ suy ra $A-M=-\dfrac{1}{2}(B-M)$, do đó $M=\dfrac{A+\dfrac{1}{2}B}{1+\dfrac{1}{2}}=(-\dfrac{1}{3};\dfrac{11}{3};-\dfrac{11}{3})$.}]{Trong không gian $Oxyz$, cho $A(-1;4;-5)$, $B(1;3;-1)$. Tìm tọa độ điểm $M$ thoả mãn $\overrightarrow{MA}=-\dfrac{1}{2}\overrightarrow{MB}$.}
  {$(-\dfrac{1}{3};\dfrac{11}{3};-\dfrac{11}{3})$.}
  {$(\dfrac{1}{3};\dfrac{11}{3};-\dfrac{11}{3})$.}{$(-\dfrac{7}{3};\dfrac{11}{3};-\dfrac{11}{3})$.}{$(-\dfrac{1}{3};\dfrac{17}{3};-\dfrac{11}{3})$.}

\cauTN[id=C2-B7-D06-TH-048,muc=TH,loigiai={Từ $\overrightarrow{MA}=3\overrightarrow{MB}$ suy ra $A-M=3(B-M)$, do đó $M=\dfrac{A-3B}{1-3}=(-\dfrac{17}{2};9;-4)$.}]{Trong không gian $Oxyz$, cho $A(5;-3;5)$, $B(-4;5;-1)$. Tìm tọa độ điểm $M$ thoả mãn $\overrightarrow{MA}=3\overrightarrow{MB}$.}
  {$(-\dfrac{17}{2};9;-4)$.}
  {$(-\dfrac{19}{2};9;-4)$.}{$(9;-\dfrac{17}{2};-4)$.}{$(-\dfrac{21}{2};9;-4)$.}

\cauDS[id=C2-B7-D06-TH-049,muc=TH,loigiai={Kiểm tra từng mệnh đề bằng công thức tọa độ. a) Sai. b) Sai. c) Sai. d) Đúng.}]{Trong không gian $Oxyz$, cho $A(2;5;-1)$, $B(1;5;-1)$, $C(-4;3;1)$. Xét tính đúng sai của các mệnh đề sau:}{%
  \y{Điểm $M$ thoả mãn $\overrightarrow{MA}=2\overrightarrow{MB}$ có tọa độ $(0;6;-1)$.}{s}
  \y{Tâm của hình bình hành $ABCD$ có tọa độ $(-1;4;2)$.}{s}
  \y{Gọi $M$, $N$ là trung điểm của $AB$, $AC$ thì $\overrightarrow{MN}=(-\dfrac{3}{2};-1;1)$.}{s}
  \y{Gọi $M$, $N$ là trung điểm của $AB$, $AC$ thì $\overrightarrow{MN}=(-\dfrac{5}{2};-1;1)$.}{d}}

\cauDS[id=C2-B7-D06-TH-050,muc=TH,loigiai={Kiểm tra từng mệnh đề bằng công thức tọa độ. a) Đúng. b) Sai. c) Đúng. d) Sai.}]{Trong không gian $Oxyz$, cho $A(-1;0;1)$, $B(-3;2;1)$, $C(3;-5;3)$. Xét tính đúng sai của các mệnh đề sau:}{%
  \y{Điểm $M$ thoả mãn $\overrightarrow{MA}=\dfrac{1}{2}\overrightarrow{MB}$ có tọa độ $(1;-2;1)$.}{d}
  \y{Tâm của hình bình hành $ABCD$ có tọa độ $(-1;\dfrac{5}{2};-2)$.}{s}
  \y{Gọi $M$, $N$ là trung điểm của $AB$, $AC$ thì $\overrightarrow{MN}=(3;-\dfrac{7}{2};1)$.}{d}
  \y{Tâm của hình bình hành $ABCD$ có tọa độ $(1;-\dfrac{1}{2};2)$.}{s}}

\cauDS[id=C2-B7-D06-TH-051,muc=TH,loigiai={Kiểm tra từng mệnh đề bằng công thức tọa độ. a) Đúng. b) Đúng. c) Sai. d) Sai.}]{Trong không gian $Oxyz$, cho $A(2;0;2)$, $B(-2;-4;5)$, $C(-1;-1;5)$, $D(-2;-2;0)$. Xét tính đúng sai của các mệnh đề sau:}{%
  \y{Trọng tâm tứ diện $ABCD$ có tọa độ $(-\dfrac{3}{4};-\dfrac{7}{4};3)$.}{d}
  \y{Tâm của hình bình hành $ABCD$ có tọa độ $(\dfrac{1}{2};-\dfrac{1}{2};\dfrac{7}{2})$.}{d}
  \y{Gọi $M$, $N$ là trung điểm của $AB$, $AC$ thì $\overrightarrow{MN}=(\dfrac{3}{2};\dfrac{1}{2};0)$.}{s}
  \y{Trọng tâm tứ diện $ABCD$ có tọa độ $(-\dfrac{7}{4};-\dfrac{7}{4};3)$.}{s}}

\cauDS[id=C2-B7-D06-TH-052,muc=TH,loigiai={Kiểm tra từng mệnh đề bằng công thức tọa độ. a) Đúng. b) Đúng. c) Đúng. d) Sai.}]{Trong không gian $Oxyz$, cho $A(4;2;-1)$, $B(-1;1;-3)$, $C(-1;1;3)$, $D(5;4;1)$. Xét tính đúng sai của các mệnh đề sau:}{%
  \y{Tâm của hình bình hành $ABCD$ có tọa độ $(\dfrac{3}{2};\dfrac{3}{2};1)$.}{d}
  \y{Điểm $M$ thoả mãn $\overrightarrow{MA}=3\overrightarrow{MB}$ có tọa độ $(-\dfrac{7}{2};\dfrac{1}{2};-4)$.}{d}
  \y{Trọng tâm tứ diện $ABCD$ có tọa độ $(\dfrac{7}{4};2;0)$.}{d}
  \y{Trọng tâm tứ diện $ABCD$ có tọa độ $(\dfrac{2}{3};\dfrac{4}{3};-\dfrac{1}{3})$.}{s}}

\cauDS[id=C2-B7-D06-TH-053,muc=TH,loigiai={Kiểm tra từng mệnh đề bằng công thức tọa độ. a) Đúng. b) Đúng. c) Sai. d) Sai.}]{Trong không gian $Oxyz$, cho $A(-1;-4;-4)$, $B(-1;-4;-2)$, $C(-2;-2;-5)$, $D(2;-5;1)$. Xét tính đúng sai của các mệnh đề sau:}{%
  \y{Điểm $M$ thoả mãn $\overrightarrow{MA}=-1\overrightarrow{MB}$ có tọa độ $(-1;-4;-3)$.}{d}
  \y{Gọi $M$, $N$ là trung điểm của $AB$, $AC$ thì $\overrightarrow{MN}=(-\dfrac{1}{2};1;-\dfrac{3}{2})$.}{d}
  \y{Trọng tâm tứ diện $ABCD$ có tọa độ $(-\dfrac{15}{4};-\dfrac{1}{2};-\dfrac{5}{2})$.}{s}
  \y{Điểm $M$ thoả mãn $\overrightarrow{MA}=3\overrightarrow{MB}$ có tọa độ $(\dfrac{1}{2};2;\dfrac{1}{2})$.}{s}}

\cauDS[id=C2-B7-D06-TH-054,muc=TH,loigiai={Kiểm tra từng mệnh đề bằng công thức tọa độ. a) Đúng. b) Sai. c) Sai. d) Sai.}]{Trong không gian $Oxyz$, cho $A(5;0;-3)$, $B(-3;-5;3)$, $C(4;-1;4)$, $D(-2;5;2)$. Xét tính đúng sai của các mệnh đề sau:}{%
  \y{Tâm của hình bình hành $ABCD$ có tọa độ $(\dfrac{9}{2};-\dfrac{1}{2};\dfrac{1}{2})$.}{d}
  \y{Trọng tâm tứ diện $ABCD$ có tọa độ $(\dfrac{3}{2};-\dfrac{1}{4};1)$.}{s}
  \y{Gọi $M$, $N$ là trung điểm của $AB$, $AC$ thì $\overrightarrow{MN}=(\dfrac{7}{2};2;\dfrac{3}{2})$.}{s}
  \y{Tâm của hình bình hành $ABCD$ có tọa độ $(\dfrac{9}{2};\dfrac{1}{2};\dfrac{1}{2})$.}{s}}

\cauDS[id=C2-B7-D06-TH-055,muc=TH,loigiai={Kiểm tra từng mệnh đề bằng công thức tọa độ. a) Sai. b) Sai. c) Đúng. d) Sai.}]{Trong không gian $Oxyz$, cho $A(1;5;-4)$, $B(3;-4;-5)$, $C(-1;0;4)$, $D(0;5;5)$. Xét tính đúng sai của các mệnh đề sau:}{%
  \y{Trọng tâm tứ diện $ABCD$ có tọa độ $(1;\dfrac{1}{3};-\dfrac{5}{3})$.}{s}
  \y{Tâm của hình bình hành $ABCD$ có tọa độ $(0;\dfrac{5}{2};2)$.}{s}
  \y{Trọng tâm tứ diện $ABCD$ có tọa độ $(\dfrac{3}{4};\dfrac{3}{2};0)$.}{d}
  \y{Gọi $M$, $N$ là trung điểm của $AB$, $AC$ thì $\overrightarrow{MN}=(0;2;\dfrac{9}{2})$.}{s}}

\cauDS[id=C2-B7-D06-TH-056,muc=TH,loigiai={Kiểm tra từng mệnh đề bằng công thức tọa độ. a) Sai. b) Đúng. c) Sai. d) Đúng.}]{Trong không gian $Oxyz$, cho $A(-5;1;-3)$, $B(2;2;-3)$, $C(-2;5;5)$, $D(3;5;-3)$. Xét tính đúng sai của các mệnh đề sau:}{%
  \y{Điểm $M$ thoả mãn $\overrightarrow{MA}=-1\overrightarrow{MB}$ có tọa độ $(\dfrac{3}{2};\dfrac{3}{2};-3)$.}{s}
  \y{Tâm của hình bình hành $ABCD$ có tọa độ $(-\dfrac{7}{2};3;1)$.}{d}
  \y{Trọng tâm tứ diện $ABCD$ có tọa độ $(-\dfrac{3}{2};\dfrac{13}{4};-1)$.}{s}
  \y{Trọng tâm tứ diện $ABCD$ có tọa độ $(-\dfrac{1}{2};\dfrac{13}{4};-1)$.}{d}}

\cauDS[id=C2-B7-D06-TH-057,muc=TH,loigiai={Kiểm tra từng mệnh đề bằng công thức tọa độ. a) Đúng. b) Sai. c) Đúng. d) Đúng.}]{Trong không gian $Oxyz$, cho $A(1;5;2)$, $B(-3;1;4)$, $C(3;-2;-5)$. Xét tính đúng sai của các mệnh đề sau:}{%
  \y{Tâm của hình bình hành $ABCD$ có tọa độ $(2;\dfrac{3}{2};-\dfrac{3}{2})$.}{d}
  \y{Gọi $M$, $N$ là trung điểm của $AB$, $AC$ thì $\overrightarrow{MN}=(3;-\dfrac{3}{2};-\dfrac{7}{2})$.}{s}
  \y{Điểm $M$ thoả mãn $\overrightarrow{MA}=-\dfrac{1}{2}\overrightarrow{MB}$ có tọa độ $(-\dfrac{1}{3};\dfrac{11}{3};\dfrac{8}{3})$.}{d}
  \y{Điểm $M$ thoả mãn $\overrightarrow{MA}=-2\overrightarrow{MB}$ có tọa độ $(-\dfrac{5}{3};\dfrac{7}{3};\dfrac{10}{3})$.}{d}}

\cauDS[id=C2-B7-D06-TH-058,muc=TH,loigiai={Kiểm tra từng mệnh đề bằng công thức tọa độ. a) Đúng. b) Sai. c) Sai. d) Sai.}]{Trong không gian $Oxyz$, cho $A(1;-3;-1)$, $B(-5;4;0)$, $C(1;0;3)$, $D(5;1;-5)$. Xét tính đúng sai của các mệnh đề sau:}{%
  \y{Điểm $M$ thoả mãn $\overrightarrow{MA}=-\dfrac{1}{2}\overrightarrow{MB}$ có tọa độ $(-1;-\dfrac{2}{3};-\dfrac{2}{3})$.}{d}
  \y{Gọi $M$, $N$ là trung điểm của $AB$, $AC$ thì $\overrightarrow{MN}=(3;0;\dfrac{3}{2})$.}{s}
  \y{Trọng tâm tứ diện $ABCD$ có tọa độ $(\dfrac{1}{2};\dfrac{1}{2};\dfrac{5}{4})$.}{s}
  \y{Điểm $M$ thoả mãn $\overrightarrow{MA}=3\overrightarrow{MB}$ có tọa độ $(-8;\dfrac{15}{2};\dfrac{5}{2})$.}{s}}

\cauDS[id=C2-B7-D06-TH-059,muc=TH,loigiai={Kiểm tra từng mệnh đề bằng công thức tọa độ. a) Sai. b) Đúng. c) Sai. d) Sai.}]{Trong không gian $Oxyz$, cho $A(0;4;-5)$, $B(5;3;3)$, $C(1;-2;3)$, $D(4;0;1)$. Xét tính đúng sai của các mệnh đề sau:}{%
  \y{Gọi $M$, $N$ là trung điểm của $AB$, $AC$ thì $\overrightarrow{MN}=(-2;0;-\dfrac{5}{2})$.}{s}
  \y{Trọng tâm tứ diện $ABCD$ có tọa độ $(\dfrac{5}{2};\dfrac{5}{4};\dfrac{1}{2})$.}{d}
  \y{Tâm của hình bình hành $ABCD$ có tọa độ $(\dfrac{1}{2};1;1)$.}{s}
  \y{Điểm $M$ thoả mãn $\overrightarrow{MA}=-2\overrightarrow{MB}$ có tọa độ $(\dfrac{10}{3};\dfrac{10}{3};\dfrac{4}{3})$.}{s}}

\cauDS[id=C2-B7-D06-TH-060,muc=TH,loigiai={Kiểm tra từng mệnh đề bằng công thức tọa độ. a) Đúng. b) Sai. c) Sai. d) Đúng.}]{Trong không gian $Oxyz$, cho $A(5;2;-1)$, $B(5;4;3)$, $C(5;2;-2)$. Xét tính đúng sai của các mệnh đề sau:}{%
  \y{Điểm $M$ thoả mãn $\overrightarrow{MA}=-\dfrac{1}{2}\overrightarrow{MB}$ có tọa độ $(5;\dfrac{8}{3};\dfrac{1}{3})$.}{d}
  \y{Gọi $M$, $N$ là trung điểm của $AB$, $AC$ thì $\overrightarrow{MN}=(0;-2;-5)$.}{s}
  \y{Gọi $M$, $N$ là trung điểm của $AB$, $AC$ thì $\overrightarrow{MN}=(0;-1;-\dfrac{1}{2})$.}{s}
  \y{Tâm của hình bình hành $ABCD$ có tọa độ $(5;2;-\dfrac{3}{2})$.}{d}}

\cauDS[id=C2-B7-D06-TH-061,muc=TH,loigiai={Kiểm tra từng mệnh đề bằng công thức tọa độ. a) Sai. b) Sai. c) Đúng. d) Đúng.}]{Trong không gian $Oxyz$, cho $A(1;2;2)$, $B(-1;4;0)$, $C(3;-2;5)$, $D(-4;5;-4)$. Xét tính đúng sai của các mệnh đề sau:}{%
  \y{Điểm $M$ thoả mãn $\overrightarrow{MA}=-2\overrightarrow{MB}$ có tọa độ $(\dfrac{1}{3};-\dfrac{10}{3};-\dfrac{2}{3})$.}{s}
  \y{Tâm của hình bình hành $ABCD$ có tọa độ $(-2;0;-\dfrac{7}{2})$.}{s}
  \y{Gọi $M$, $N$ là trung điểm của $AB$, $AC$ thì $\overrightarrow{MN}=(2;-3;\dfrac{5}{2})$.}{d}
  \y{Trọng tâm tứ diện $ABCD$ có tọa độ $(-\dfrac{1}{4};\dfrac{9}{4};\dfrac{3}{4})$.}{d}}

\cauDS[id=C2-B7-D06-TH-062,muc=TH,loigiai={Kiểm tra từng mệnh đề bằng công thức tọa độ. a) Sai. b) Sai. c) Đúng. d) Đúng.}]{Trong không gian $Oxyz$, cho $A(-4;2;5)$, $B(4;3;3)$, $C(-3;-1;5)$. Xét tính đúng sai của các mệnh đề sau:}{%
  \y{Gọi $M$, $N$ là trung điểm của $AB$, $AC$ thì $\overrightarrow{MN}=(-\dfrac{7}{2};2;1)$.}{s}
  \y{Tâm của hình bình hành $ABCD$ có tọa độ $(-\dfrac{3}{2};\dfrac{1}{2};5)$.}{s}
  \y{Tâm của hình bình hành $ABCD$ có tọa độ $(-\dfrac{7}{2};\dfrac{1}{2};5)$.}{d}
  \y{Điểm $M$ thoả mãn $\overrightarrow{MA}=-\dfrac{1}{2}\overrightarrow{MB}$ có tọa độ $(-\dfrac{4}{3};\dfrac{7}{3};\dfrac{13}{3})$.}{d}}

\cauDS[id=C2-B7-D06-TH-063,muc=TH,loigiai={Kiểm tra từng mệnh đề bằng công thức tọa độ. a) Sai. b) Sai. c) Đúng. d) Sai.}]{Trong không gian $Oxyz$, cho $A(4;4;5)$, $B(0;-2;2)$, $C(2;-4;-3)$, $D(2;-1;-4)$. Xét tính đúng sai của các mệnh đề sau:}{%
  \y{Gọi $M$, $N$ là trung điểm của $AB$, $AC$ thì $\overrightarrow{MN}=(1;-\dfrac{5}{2};-1)$.}{s}
  \y{Tâm của hình bình hành $ABCD$ có tọa độ $(3;2;1)$.}{s}
  \y{Trọng tâm tứ diện $ABCD$ có tọa độ $(2;-\dfrac{3}{4};0)$.}{d}
  \y{Gọi $M$, $N$ là trung điểm của $AB$, $AC$ thì $\overrightarrow{MN}=(-1;1;-\dfrac{5}{2})$.}{s}}

\cauDS[id=C2-B7-D06-TH-064,muc=TH,loigiai={Kiểm tra từng mệnh đề bằng công thức tọa độ. a) Sai. b) Đúng. c) Sai. d) Đúng.}]{Trong không gian $Oxyz$, cho $A(4;-3;3)$, $B(-1;3;1)$, $C(5;-4;2)$, $D(3;1;-4)$. Xét tính đúng sai của các mệnh đề sau:}{%
  \y{Điểm $M$ thoả mãn $\overrightarrow{MA}=-\dfrac{1}{2}\overrightarrow{MB}$ có tọa độ $(\dfrac{7}{3};-1;\dfrac{10}{3})$.}{s}
  \y{Điểm $M$ thoả mãn $\overrightarrow{MA}=\dfrac{1}{2}\overrightarrow{MB}$ có tọa độ $(9;-9;5)$.}{d}
  \y{Trọng tâm tứ diện $ABCD$ có tọa độ $(\dfrac{11}{4};\dfrac{3}{4};\dfrac{1}{2})$.}{s}
  \y{Tâm của hình bình hành $ABCD$ có tọa độ $(\dfrac{9}{2};-\dfrac{7}{2};\dfrac{5}{2})$.}{d}}

\cauTLN[id=C2-B7-D06-TH-065,muc=TH,loigiai={Vì $ABCD$ là hình bình hành nên $D=A+C-B=(9;-4;-2)$. Tâm $I$ của hình bình hành là trung điểm của $AC$ (cũng là trung điểm $BD$): $I=(3;0;1)$.}]{Trong không gian $Oxyz$, cho $A(4;-3;-1)$, $B(-3;4;4)$, $C(2;3;3)$. Gọi $I$ là tâm của hình bình hành $ABCD$. Tính $x_I+y_I+z_I$.}{4}

\cauTLN[id=C2-B7-D06-TH-066,muc=TH,loigiai={Vì $ABCD$ là hình bình hành nên $D=A+C-B=(7;5;-1)$. Tâm $I$ của hình bình hành là trung điểm của $AC$ (cũng là trung điểm $BD$): $I=(2;2;1)$.}]{Trong không gian $Oxyz$, cho $A(-1;0;4)$, $B(-3;-1;3)$, $C(5;4;-2)$. Gọi $I$ là tâm của hình bình hành $ABCD$. Tính $x_I+y_I+z_I$.}{5}

\cauTLN[id=C2-B7-D06-TH-067,muc=TH,loigiai={Vì $ABCD$ là hình bình hành nên $D=A+C-B=(-3;-4;-1)$. Tâm $I$ của hình bình hành là trung điểm của $AC$ (cũng là trung điểm $BD$): $I=(1;-3;-2)$.}]{Trong không gian $Oxyz$, cho $A(-3;-3;-2)$, $B(5;-2;-3)$, $C(5;-3;-2)$. Gọi $I$ là tâm của hình bình hành $ABCD$. Tính $x_I+y_I+z_I$.}{-4}

\cauTLN[id=C2-B7-D06-TH-068,muc=TH,loigiai={Vì $ABCD$ là hình bình hành nên $D=A+C-B=(-7;7;-2)$. Tâm $I$ của hình bình hành là trung điểm của $AC$ (cũng là trung điểm $BD$): $I=(-4;1;-2)$.}]{Trong không gian $Oxyz$, cho $A(-4;0;-3)$, $B(-1;-5;-2)$, $C(-4;2;-1)$. Gọi $I$ là tâm của hình bình hành $ABCD$. Tính $x_I+y_I+z_I$.}{-5}

\cauTLN[id=C2-B7-D06-TH-069,muc=TH,loigiai={Vì $ABCD$ là hình bình hành nên $D=A+C-B=(5;4;-2)$. Tâm $I$ của hình bình hành là trung điểm của $AC$ (cũng là trung điểm $BD$): $I=(3;4;-1)$.}]{Trong không gian $Oxyz$, cho $A(3;4;-3)$, $B(1;4;0)$, $C(3;4;1)$. Gọi $I$ là tâm của hình bình hành $ABCD$. Tính $x_I+y_I+z_I$.}{6}

\cauTLN[id=C2-B7-D06-TH-070,muc=TH,loigiai={Trọng tâm tứ diện $ABCD$ có tọa độ là trung bình cộng tọa độ bốn đỉnh: $G=(-1;0;0)$.}]{Trong không gian $Oxyz$, cho $A(4;-4;2)$, $B(-3;4;-2)$, $C(-1;5;-4)$, $D(-4;-5;4)$. Gọi $G$ là trọng tâm tứ diện $ABCD$. Tính $x_G+y_G+z_G$.}{-1}

\cauTLN[id=C2-B7-D06-TH-071,muc=TH,loigiai={Trọng tâm tứ diện $ABCD$ có tọa độ là trung bình cộng tọa độ bốn đỉnh: $G=(0;-2;1)$.}]{Trong không gian $Oxyz$, cho $A(5;0;5)$, $B(-4;-2;-2)$, $C(-2;-3;-4)$, $D(1;-3;5)$. Gọi $G$ là trọng tâm tứ diện $ABCD$. Tính $x_G+y_G+z_G$.}{-1}

\cauTLN[id=C2-B7-D06-TH-072,muc=TH,loigiai={Trọng tâm tứ diện $ABCD$ có tọa độ là trung bình cộng tọa độ bốn đỉnh: $G=(3;-1;0)$.}]{Trong không gian $Oxyz$, cho $A(3;-2;5)$, $B(5;3;-3)$, $C(2;-5;-4)$, $D(2;0;2)$. Gọi $G$ là trọng tâm tứ diện $ABCD$. Tính $x_G+y_G+z_G$.}{2}

\cauTLN[id=C2-B7-D06-TH-073,muc=TH,loigiai={Trọng tâm tứ diện $ABCD$ có tọa độ là trung bình cộng tọa độ bốn đỉnh: $G=(0;-1;1)$.}]{Trong không gian $Oxyz$, cho $A(5;1;2)$, $B(-2;0;2)$, $C(-5;0;1)$, $D(2;-5;-1)$. Gọi $G$ là trọng tâm tứ diện $ABCD$. Tính $x_G+y_G+z_G$.}{0}

\cauTLN[id=C2-B7-D06-TH-074,muc=TH,loigiai={Trọng tâm tứ diện $ABCD$ có tọa độ là trung bình cộng tọa độ bốn đỉnh: $G=(-3;-2;-2)$.}]{Trong không gian $Oxyz$, cho $A(-4;-5;-5)$, $B(-3;2;1)$, $C(0;-2;-4)$, $D(-5;-3;0)$. Gọi $G$ là trọng tâm tứ diện $ABCD$. Tính $x_G+y_G+z_G$.}{-7}

\cauTLN[id=C2-B7-D06-TH-075,muc=TH,loigiai={Gọi $M$, $N$ lần lượt là trung điểm $AB$, $AC$: $M=(-\dfrac{1}{2};1;-\dfrac{1}{2})$, $N=(-\dfrac{3}{2};4;-\dfrac{1}{2})$ nên $\overrightarrow{MN}=(-1;3;0)$.}]{Trong không gian $Oxyz$, cho $A(-2;3;-5)$, $B(1;-1;4)$, $C(-1;5;4)$. Gọi $M$, $N$ là trung điểm của $AB$ và $AC$. Tính $x+y+z$ với $(x;y;z)$ là tọa độ vectơ $\overrightarrow{MN}$.}{2}

\cauTLN[id=C2-B7-D06-TH-076,muc=TH,loigiai={Gọi $M$, $N$ lần lượt là trung điểm $AB$, $AC$: $M=(\dfrac{3}{2};0;-\dfrac{7}{2})$, $N=(\dfrac{3}{2};-3;-\dfrac{3}{2})$ nên $\overrightarrow{MN}=(0;-3;2)$.}]{Trong không gian $Oxyz$, cho $A(0;-2;-3)$, $B(3;2;-4)$, $C(3;-4;0)$. Gọi $M$, $N$ là trung điểm của $AB$ và $AC$. Tính $x+y+z$ với $(x;y;z)$ là tọa độ vectơ $\overrightarrow{MN}$.}{-1}

\cauTLN[id=C2-B7-D06-TH-077,muc=TH,loigiai={Gọi $M$, $N$ lần lượt là trung điểm $AB$, $AC$: $M=(0;\dfrac{3}{2};4)$, $N=(-4;\dfrac{1}{2};0)$ nên $\overrightarrow{MN}=(-4;-1;-4)$.}]{Trong không gian $Oxyz$, cho $A(-3;4;3)$, $B(3;-1;5)$, $C(-5;-3;-3)$. Gọi $M$, $N$ là trung điểm của $AB$ và $AC$. Tính $x+y+z$ với $(x;y;z)$ là tọa độ vectơ $\overrightarrow{MN}$.}{-9}

\cauTLN[id=C2-B7-D06-TH-078,muc=TH,loigiai={Gọi $M$, $N$ lần lượt là trung điểm $AB$, $AC$: $M=(\dfrac{1}{2};-\dfrac{1}{2};\dfrac{3}{2})$, $N=(-\dfrac{1}{2};\dfrac{7}{2};-\dfrac{1}{2})$ nên $\overrightarrow{MN}=(-1;4;-2)$.}]{Trong không gian $Oxyz$, cho $A(-3;2;0)$, $B(4;-3;3)$, $C(2;5;-1)$. Gọi $M$, $N$ là trung điểm của $AB$ và $AC$. Tính $x+y+z$ với $(x;y;z)$ là tọa độ vectơ $\overrightarrow{MN}$.}{1}

\cauTLN[id=C2-B7-D06-TH-079,muc=TH,loigiai={Gọi $M$, $N$ lần lượt là trung điểm $AB$, $AC$: $M=(1;2;\dfrac{5}{2})$, $N=(3;-1;-\dfrac{1}{2})$ nên $\overrightarrow{MN}=(2;-3;-3)$.}]{Trong không gian $Oxyz$, cho $A(4;1;0)$, $B(-2;3;5)$, $C(2;-3;-1)$. Gọi $M$, $N$ là trung điểm của $AB$ và $AC$. Tính $x+y+z$ với $(x;y;z)$ là tọa độ vectơ $\overrightarrow{MN}$.}{-4}

\cauTLN[id=C2-B7-D06-TH-080,muc=TH,loigiai={Gọi $M$, $N$ lần lượt là trung điểm $AB$, $AC$: $M=(3;0;5)$, $N=(2;-2;0)$ nên $\overrightarrow{MN}=(-1;-2;-5)$.}]{Trong không gian $Oxyz$, cho $A(1;-3;5)$, $B(5;3;5)$, $C(3;-1;-5)$. Gọi $M$, $N$ là trung điểm của $AB$ và $AC$. Tính $x+y+z$ với $(x;y;z)$ là tọa độ vectơ $\overrightarrow{MN}$.}{-8}

\cauTN[id=C2-B7-D06-VD-001,muc=VD,loigiai={Vì $M,N,P$ là trung điểm $BC,CA,AB$ nên $APMN$ là hình bình hành: $\overrightarrow{AP}=\overrightarrow{NM}$ hay $A=N+P-M=(0;-4;4)$; tương tự $B=M+P-N=(-4;0;-6)$, $C=M+N-P=(4;2;2)$.}]{Tam giác $ABC$ có các trung điểm của $BC$, $CA$, $AB$ lần lượt là $M(0;1;-2)$, $N(2;-1;3)$, $P(-2;-2;-1)$. Tọa độ đỉnh $A$ là}
  {$(0;-4;4)$.}
  {$(0;4;-4)$.}{$(0;-5;4)$.}{$(-1;-4;4)$.}

\cauTN[id=C2-B7-D06-VD-002,muc=VD,loigiai={Vì $M,N,P$ là trung điểm $BC,CA,AB$ nên $APMN$ là hình bình hành: $\overrightarrow{AP}=\overrightarrow{NM}$ hay $A=N+P-M=(-2;4;-5)$; tương tự $B=M+P-N=(2;4;-3)$, $C=M+N-P=(4;-2;1)$.}]{Tam giác $ABC$ có các trung điểm của $BC$, $CA$, $AB$ lần lượt là $M(3;1;-1)$, $N(1;1;-2)$, $P(0;4;-4)$. Tọa độ đỉnh $C$ là}
  {$(4;-2;1)$.}
  {$(4;-2;-1)$.}{$(4;-4;1)$.}{$(4;-1;1)$.}

\cauTN[id=C2-B7-D06-VD-003,muc=VD,loigiai={Vì $M,N,P$ là trung điểm $BC,CA,AB$ nên $APMN$ là hình bình hành: $\overrightarrow{AP}=\overrightarrow{NM}$ hay $A=N+P-M=(4;-1;2)$; tương tự $B=M+P-N=(-6;3;-6)$, $C=M+N-P=(-2;3;4)$.}]{Tam giác $ABC$ có các trung điểm của $BC$, $CA$, $AB$ lần lượt là $M(-4;3;-1)$, $N(1;1;3)$, $P(-1;1;-2)$. Tọa độ đỉnh $A$ là}
  {$(4;-1;2)$.}
  {$(5;-1;2)$.}{$(-4;1;-2)$.}{$(4;1;2)$.}

\cauTN[id=C2-B7-D06-VD-004,muc=VD,loigiai={Vì $M,N,P$ là trung điểm $BC,CA,AB$ nên $APMN$ là hình bình hành: $\overrightarrow{AP}=\overrightarrow{NM}$ hay $A=N+P-M=(2;-2;-2)$; tương tự $B=M+P-N=(-6;-4;0)$, $C=M+N-P=(-6;-4;2)$.}]{Tam giác $ABC$ có các trung điểm của $BC$, $CA$, $AB$ lần lượt là $M(-6;-4;1)$, $N(-2;-3;0)$, $P(-2;-3;-1)$. Tọa độ đỉnh $A$ là}
  {$(2;-2;-2)$.}
  {$(8;2;-3)$.}{$(2;-2;-3)$.}{$(2;-2;2)$.}

\cauTN[id=C2-B7-D06-VD-005,muc=VD,loigiai={Vì $M,N,P$ là trung điểm $BC,CA,AB$ nên $APMN$ là hình bình hành: $\overrightarrow{AP}=\overrightarrow{NM}$ hay $A=N+P-M=(-3;-1;-6)$; tương tự $B=M+P-N=(-1;-3;4)$, $C=M+N-P=(-1;5;4)$.}]{Tam giác $ABC$ có các trung điểm của $BC$, $CA$, $AB$ lần lượt là $M(-1;1;4)$, $N(-2;2;-1)$, $P(-2;-2;-1)$. Tọa độ đỉnh $B$ là}
  {$(-1;-3;4)$.}
  {$(-1;-2;4)$.}{$(-1;-3;5)$.}{$(1;-3;4)$.}

\cauTN[id=C2-B7-D06-VD-006,muc=VD,loigiai={Vì $M,N,P$ là trung điểm $BC,CA,AB$ nên $APMN$ là hình bình hành: $\overrightarrow{AP}=\overrightarrow{NM}$ hay $A=N+P-M=(5;-3;3)$; tương tự $B=M+P-N=(-1;-3;1)$, $C=M+N-P=(-1;-1;5)$.}]{Tam giác $ABC$ có các trung điểm của $BC$, $CA$, $AB$ lần lượt là $M(-1;-2;3)$, $N(2;-2;4)$, $P(2;-3;2)$. Tọa độ đỉnh $C$ là}
  {$(-1;-1;5)$.}
  {$(1;-1;5)$.}{$(-1;-1;7)$.}{$(-2;-1;5)$.}

\cauTN[id=C2-B7-D06-VD-007,muc=VD,loigiai={Vì $M,N,P$ là trung điểm $BC,CA,AB$ nên $APMN$ là hình bình hành: $\overrightarrow{AP}=\overrightarrow{NM}$ hay $A=N+P-M=(-5;2;-5)$; tương tự $B=M+P-N=(3;0;-5)$, $C=M+N-P=(-1;-6;-3)$.}]{Tam giác $ABC$ có các trung điểm của $BC$, $CA$, $AB$ lần lượt là $M(1;-3;-4)$, $N(-3;-2;-4)$, $P(-1;1;-5)$. Tọa độ đỉnh $C$ là}
  {$(-1;-6;-3)$.}
  {$(0;-6;-3)$.}{$(-1;-4;-3)$.}{$(-1;-6;-5)$.}

\cauTN[id=C2-B7-D06-VD-008,muc=VD,loigiai={Vì $M,N,P$ là trung điểm $BC,CA,AB$ nên $APMN$ là hình bình hành: $\overrightarrow{AP}=\overrightarrow{NM}$ hay $A=N+P-M=(3;-5;-2)$; tương tự $B=M+P-N=(-5;1;-6)$, $C=M+N-P=(-3;-5;-2)$.}]{Tam giác $ABC$ có các trung điểm của $BC$, $CA$, $AB$ lần lượt là $M(-4;-2;-4)$, $N(0;-5;-2)$, $P(-1;-2;-4)$. Tọa độ đỉnh $C$ là}
  {$(-3;-5;-2)$.}
  {$(-2;-5;-3)$.}{$(-2;-5;-2)$.}{$(-3;-5;-1)$.}

\cauTN[id=C2-B7-D06-VD-009,muc=VD,loigiai={Vì $M,N,P$ là trung điểm $BC,CA,AB$ nên $APMN$ là hình bình hành: $\overrightarrow{AP}=\overrightarrow{NM}$ hay $A=N+P-M=(3;3;2)$; tương tự $B=M+P-N=(-1;-3;-2)$, $C=M+N-P=(-1;5;0)$.}]{Tam giác $ABC$ có các trung điểm của $BC$, $CA$, $AB$ lần lượt là $M(-1;1;-1)$, $N(1;4;1)$, $P(1;0;0)$. Tọa độ đỉnh $B$ là}
  {$(-1;-3;-2)$.}
  {$(-1;-2;-2)$.}{$(-1;3;-2)$.}{$(-2;-3;-1)$.}

\cauTN[id=C2-B7-D06-VD-010,muc=VD,loigiai={Vì $M,N,P$ là trung điểm $BC,CA,AB$ nên $APMN$ là hình bình hành: $\overrightarrow{AP}=\overrightarrow{NM}$ hay $A=N+P-M=(-5;-4;-3)$; tương tự $B=M+P-N=(-5;0;-1)$, $C=M+N-P=(-1;0;-1)$.}]{Tam giác $ABC$ có các trung điểm của $BC$, $CA$, $AB$ lần lượt là $M(-3;0;-1)$, $N(-3;-2;-2)$, $P(-5;-2;-2)$. Tọa độ đỉnh $C$ là}
  {$(-1;0;-1)$.}
  {$(-1;-1;0)$.}{$(1;0;-1)$.}{$(-1;2;-1)$.}

\cauTN[id=C2-B7-D06-VD-011,muc=VD,loigiai={Vì $M,N,P$ là trung điểm $BC,CA,AB$ nên $APMN$ là hình bình hành: $\overrightarrow{AP}=\overrightarrow{NM}$ hay $A=N+P-M=(-5;-2;2)$; tương tự $B=M+P-N=(-1;-4;-6)$, $C=M+N-P=(-1;0;0)$.}]{Tam giác $ABC$ có các trung điểm của $BC$, $CA$, $AB$ lần lượt là $M(-1;-2;-3)$, $N(-3;-1;1)$, $P(-3;-3;-2)$. Tọa độ đỉnh $A$ là}
  {$(-5;-2;2)$.}
  {$(-5;-2;3)$.}{$(-5;-2;1)$.}{$(-5;-2;4)$.}

\cauTN[id=C2-B7-D06-VD-012,muc=VD,loigiai={Vì $M,N,P$ là trung điểm $BC,CA,AB$ nên $APMN$ là hình bình hành: $\overrightarrow{AP}=\overrightarrow{NM}$ hay $A=N+P-M=(-4;2;-3)$; tương tự $B=M+P-N=(-6;0;1)$, $C=M+N-P=(-2;0;-5)$.}]{Tam giác $ABC$ có các trung điểm của $BC$, $CA$, $AB$ lần lượt là $M(-4;0;-2)$, $N(-3;1;-4)$, $P(-5;1;-1)$. Tọa độ đỉnh $C$ là}
  {$(-2;0;-5)$.}
  {$(-5;0;-2)$.}{$(0;-2;-5)$.}{$(-2;1;-5)$.}

\cauTN[id=C2-B7-D06-VD-013,muc=VD,loigiai={Tâm $I$ là trung điểm của $AC$ và $BD$ nên $D=2I-B=(-6;1;1)$.}]{Hình bình hành $ABCD$ có tâm $I(-2;-\dfrac{1}{2};2)$ và hai đỉnh $A(-4;2;1)$, $B(2;-2;3)$. Tọa độ đỉnh $D$ là}
  {$(-6;1;1)$.}
  {$(-4;1;1)$.}{$(-6;1;3)$.}{$(-6;3;1)$.}

\cauTN[id=C2-B7-D06-VD-014,muc=VD,loigiai={Tâm $I$ là trung điểm của $AC$ và $BD$ nên $D=2I-B=(5;1;3)$.}]{Hình bình hành $ABCD$ có tâm $I(\dfrac{5}{2};-1;\dfrac{5}{2})$ và hai đỉnh $A(1;-2;2)$, $B(0;-3;2)$. Tọa độ đỉnh $D$ là}
  {$(5;1;3)$.}
  {$(-5;1;3)$.}{$(5;1;5)$.}{$(7;1;3)$.}

\cauTN[id=C2-B7-D06-VD-015,muc=VD,loigiai={Tâm $I$ là trung điểm của $AC$ và $BD$ nên $D=2I-B=(-1;3;4)$.}]{Hình bình hành $ABCD$ có tâm $I(0;\dfrac{5}{2};2)$ và hai đỉnh $A(3;3;2)$, $B(1;2;0)$. Tọa độ đỉnh $D$ là}
  {$(-1;3;4)$.}
  {$(1;3;4)$.}{$(-1;4;3)$.}{$(4;5;2)$.}

\cauTN[id=C2-B7-D06-VD-016,muc=VD,loigiai={Tâm $I$ là trung điểm của $AC$ và $BD$ nên $D=2I-B=(-4;4;6)$.}]{Hình bình hành $ABCD$ có tâm $I(-\dfrac{1}{2};3;1)$ và hai đỉnh $A(2;2;1)$, $B(3;2;-4)$. Tọa độ đỉnh $D$ là}
  {$(-4;4;6)$.}
  {$(-4;4;8)$.}{$(-4;4;7)$.}{$(\dfrac{7}{2};-1;-5)$.}

\cauTN[id=C2-B7-D06-VD-017,muc=VD,loigiai={Tâm $I$ là trung điểm của $AC$ và $BD$ nên $D=2I-B=(0;2;-7)$.}]{Hình bình hành $ABCD$ có tâm $I(\dfrac{1}{2};\dfrac{5}{2};-3)$ và hai đỉnh $A(4;2;-4)$, $B(1;3;1)$. Tọa độ đỉnh $D$ là}
  {$(0;2;-7)$.}
  {$(0;2;-6)$.}{$(5;5;-3)$.}{$(0;2;7)$.}

\cauTN[id=C2-B7-D06-VD-018,muc=VD,loigiai={Tâm $I$ là trung điểm của $AC$ và $BD$ nên $D=2I-B=(-1;6;1)$.}]{Hình bình hành $ABCD$ có tâm $I(0;2;-\dfrac{1}{2})$ và hai đỉnh $A(-4;3;0)$, $B(1;-2;-2)$. Tọa độ đỉnh $D$ là}
  {$(-1;6;1)$.}
  {$(6;-1;1)$.}{$(-1;8;1)$.}{$(-1;6;0)$.}

\cauTN[id=C2-B7-D06-VD-019,muc=VD,loigiai={Tâm $I$ là trung điểm của $AC$ và $BD$ nên $D=2I-B=(-7;0;-5)$.}]{Hình bình hành $ABCD$ có tâm $I(-\dfrac{3}{2};-2;-\dfrac{7}{2})$ và hai đỉnh $A(-1;0;-4)$, $B(4;-4;-2)$. Tọa độ đỉnh $D$ là}
  {$(-7;0;-5)$.}
  {$(3;-4;-6)$.}{$(-8;0;-5)$.}{$(-7;-1;-5)$.}

\cauTN[id=C2-B7-D06-VD-020,muc=VD,loigiai={Tâm $I$ là trung điểm của $AC$ và $BD$ nên $D=2I-B=(7;0;-1)$.}]{Hình bình hành $ABCD$ có tâm $I(2;\dfrac{3}{2};-\dfrac{5}{2})$ và hai đỉnh $A(2;2;-1)$, $B(-3;3;-4)$. Tọa độ đỉnh $D$ là}
  {$(7;0;-1)$.}
  {$(7;0;0)$.}{$(-5;\dfrac{3}{2};-\dfrac{3}{2})$.}{$(-7;0;1)$.}

\cauTN[id=C2-B7-D06-VD-021,muc=VD,loigiai={Tâm $I$ là trung điểm của $AC$ và $BD$ nên $D=2I-B=(-6;-2;-3)$.}]{Hình bình hành $ABCD$ có tâm $I(-1;0;0)$ và hai đỉnh $A(-3;4;3)$, $B(4;2;3)$. Tọa độ đỉnh $D$ là}
  {$(-6;-2;-3)$.}
  {$(-6;-2;-1)$.}{$(-6;-4;-3)$.}{$(-6;2;-3)$.}

\cauTN[id=C2-B7-D06-VD-022,muc=VD,loigiai={Tâm $I$ là trung điểm của $AC$ và $BD$ nên $D=2I-B=(-1;2;4)$.}]{Hình bình hành $ABCD$ có tâm $I(-\dfrac{5}{2};-1;\dfrac{1}{2})$ và hai đỉnh $A(-4;0;-2)$, $B(-4;-4;-3)$. Tọa độ đỉnh $D$ là}
  {$(-1;2;4)$.}
  {$(1;2;4)$.}{$(-2;2;4)$.}{$(0;2;4)$.}

\cauTN[id=C2-B7-D06-VD-023,muc=VD,loigiai={Tâm $I$ là trung điểm của $AC$ và $BD$ nên $D=2I-B=(6;-3;0)$.}]{Hình bình hành $ABCD$ có tâm $I(1;-1;2)$ và hai đỉnh $A(4;-1;0)$, $B(-4;1;4)$. Tọa độ đỉnh $D$ là}
  {$(6;-3;0)$.}
  {$(6;-3;-1)$.}{$(6;-1;0)$.}{$(6;3;0)$.}

\cauTN[id=C2-B7-D06-VD-024,muc=VD,loigiai={Tâm $I$ là trung điểm của $AC$ và $BD$ nên $D=2I-B=(4;4;1)$.}]{Hình bình hành $ABCD$ có tâm $I(0;\dfrac{5}{2};\dfrac{3}{2})$ và hai đỉnh $A(3;3;3)$, $B(-4;1;2)$. Tọa độ đỉnh $D$ là}
  {$(4;4;1)$.}
  {$(3;4;1)$.}{$(-4;4;1)$.}{$(4;4;3)$.}

\cauTN[id=C2-B7-D06-VD-025,muc=VD,loigiai={Vì $G$ là trọng tâm tứ diện nên $A+B+C+D=4G$, suy ra $D=4G-A-B-C=(3;0;3)$.}]{Tứ diện $ABCD$ có $A(3;1;-4)$, $B(3;1;4)$, $C(-4;-4;2)$ và trọng tâm $G(\dfrac{5}{4};-\dfrac{1}{2};\dfrac{5}{4})$. Tọa độ đỉnh $D$ là}
  {$(3;0;3)$.}
  {$(1;0;3)$.}{$(5;0;3)$.}{$(3;-2;3)$.}

\cauTN[id=C2-B7-D06-VD-026,muc=VD,loigiai={Vì $G$ là trọng tâm tứ diện nên $A+B+C+D=4G$, suy ra $D=4G-A-B-C=(0;-1;1)$.}]{Tứ diện $ABCD$ có $A(1;-1;4)$, $B(0;2;-1)$, $C(1;-3;-1)$ và trọng tâm $G(\dfrac{1}{2};-\dfrac{3}{4};\dfrac{3}{4})$. Tọa độ đỉnh $D$ là}
  {$(0;-1;1)$.}
  {$(2;-1;1)$.}{$(0;1;1)$.}{$(-\dfrac{3}{2};\dfrac{5}{4};-\dfrac{5}{4})$.}

\cauTN[id=C2-B7-D06-VD-027,muc=VD,loigiai={Vì $G$ là trọng tâm tứ diện nên $A+B+C+D=4G$, suy ra $D=4G-A-B-C=(1;2;2)$.}]{Tứ diện $ABCD$ có $A(-2;-2;1)$, $B(0;1;-3)$, $C(-1;2;3)$ và trọng tâm $G(-\dfrac{1}{2};\dfrac{3}{4};\dfrac{3}{4})$. Tọa độ đỉnh $D$ là}
  {$(1;2;2)$.}
  {$(0;2;2)$.}{$(1;2;1)$.}{$(-1;-2;-2)$.}

\cauTN[id=C2-B7-D06-VD-028,muc=VD,loigiai={Vì $G$ là trọng tâm tứ diện nên $A+B+C+D=4G$, suy ra $D=4G-A-B-C=(-4;0;1)$.}]{Tứ diện $ABCD$ có $A(1;2;1)$, $B(1;-4;0)$, $C(-2;3;0)$ và trọng tâm $G(-1;\dfrac{1}{4};\dfrac{1}{2})$. Tọa độ đỉnh $D$ là}
  {$(-4;0;1)$.}
  {$(-5;0;1)$.}{$(-2;0;1)$.}{$(4;0;1)$.}

\cauTN[id=C2-B7-D06-VD-029,muc=VD,loigiai={Vì $G$ là trọng tâm tứ diện nên $A+B+C+D=4G$, suy ra $D=4G-A-B-C=(1;3;-2)$.}]{Tứ diện $ABCD$ có $A(1;1;2)$, $B(-1;4;-2)$, $C(-2;-4;-1)$ và trọng tâm $G(-\dfrac{1}{4};1;-\dfrac{3}{4})$. Tọa độ đỉnh $D$ là}
  {$(1;3;-2)$.}
  {$(1;5;-2)$.}{$(3;3;-2)$.}{$(2;3;-2)$.}

\cauTN[id=C2-B7-D06-VD-030,muc=VD,loigiai={Vì $G$ là trọng tâm tứ diện nên $A+B+C+D=4G$, suy ra $D=4G-A-B-C=(-2;-1;0)$.}]{Tứ diện $ABCD$ có $A(-1;-4;-4)$, $B(-3;1;2)$, $C(-1;2;2)$ và trọng tâm $G(-\dfrac{7}{4};-\dfrac{1}{2};0)$. Tọa độ đỉnh $D$ là}
  {$(-2;-1;0)$.}
  {$(2;-1;0)$.}{$(\dfrac{13}{4};\dfrac{1}{2};0)$.}{$(-2;-1;-2)$.}

\cauTN[id=C2-B7-D06-VD-031,muc=VD,loigiai={Vì $G$ là trọng tâm tứ diện nên $A+B+C+D=4G$, suy ra $D=4G-A-B-C=(-3;-3;-4)$.}]{Tứ diện $ABCD$ có $A(4;0;0)$, $B(2;1;-4)$, $C(-1;-1;1)$ và trọng tâm $G(\dfrac{1}{2};-\dfrac{3}{4};-\dfrac{7}{4})$. Tọa độ đỉnh $D$ là}
  {$(-3;-3;-4)$.}
  {$(-3;-1;-4)$.}{$(-3;-4;-3)$.}{$(3;3;4)$.}

\cauTN[id=C2-B7-D06-VD-032,muc=VD,loigiai={Vì $G$ là trọng tâm tứ diện nên $A+B+C+D=4G$, suy ra $D=4G-A-B-C=(-3;0;1)$.}]{Tứ diện $ABCD$ có $A(-2;-2;-1)$, $B(-3;3;4)$, $C(-4;2;-1)$ và trọng tâm $G(-3;\dfrac{3}{4};\dfrac{3}{4})$. Tọa độ đỉnh $D$ là}
  {$(-3;0;1)$.}
  {$(-3;0;2)$.}{$(6;-\dfrac{9}{4};-\dfrac{5}{4})$.}{$(-3;1;1)$.}

\cauTN[id=C2-B7-D06-VD-033,muc=VD,loigiai={Vì $G$ là trọng tâm tứ diện nên $A+B+C+D=4G$, suy ra $D=4G-A-B-C=(-3;3;-3)$.}]{Tứ diện $ABCD$ có $A(1;0;4)$, $B(-1;-4;2)$, $C(0;2;-4)$ và trọng tâm $G(-\dfrac{3}{4};\dfrac{1}{4};-\dfrac{1}{4})$. Tọa độ đỉnh $D$ là}
  {$(-3;3;-3)$.}
  {$(3;-3;-3)$.}{$(-3;3;-5)$.}{$(-3;3;3)$.}

\cauTN[id=C2-B7-D06-VD-034,muc=VD,loigiai={Vì $G$ là trọng tâm tứ diện nên $A+B+C+D=4G$, suy ra $D=4G-A-B-C=(-3;0;-3)$.}]{Tứ diện $ABCD$ có $A(-2;-4;3)$, $B(-4;4;1)$, $C(4;2;-1)$ và trọng tâm $G(-\dfrac{5}{4};\dfrac{1}{2};0)$. Tọa độ đỉnh $D$ là}
  {$(-3;0;-3)$.}
  {$(-3;2;-3)$.}{$(-3;0;3)$.}{$(\dfrac{3}{4};-\dfrac{3}{2};-3)$.}

\cauTN[id=C2-B7-D06-VD-035,muc=VD,loigiai={Vì $G$ là trọng tâm tứ diện nên $A+B+C+D=4G$, suy ra $D=4G-A-B-C=(-3;-2;1)$.}]{Tứ diện $ABCD$ có $A(4;1;2)$, $B(-2;-3;3)$, $C(1;1;0)$ và trọng tâm $G(0;-\dfrac{3}{4};\dfrac{3}{2})$. Tọa độ đỉnh $D$ là}
  {$(-3;-2;1)$.}
  {$(-3;1;-2)$.}{$(-2;-3;1)$.}{$(-3;-2;3)$.}

\cauTN[id=C2-B7-D06-VD-036,muc=VD,loigiai={Vì $G$ là trọng tâm tứ diện nên $A+B+C+D=4G$, suy ra $D=4G-A-B-C=(-3;2;4)$.}]{Tứ diện $ABCD$ có $A(-4;-2;-1)$, $B(-4;4;3)$, $C(-3;0;-1)$ và trọng tâm $G(-\dfrac{7}{2};1;\dfrac{5}{4})$. Tọa độ đỉnh $D$ là}
  {$(-3;2;4)$.}
  {$(-3;-2;4)$.}{$(-3;2;-4)$.}{$(3;2;4)$.}

\cauDS[id=C2-B7-D06-VD-037,muc=VD,loigiai={Kiểm tra từng mệnh đề bằng công thức tọa độ. a) Đúng. b) Đúng. c) Sai. d) Sai.}]{Trong không gian $Oxyz$, cho $A(-3;0;1)$, $B(-1;0;-1)$, $C(1;4;4)$. Xét tính đúng sai của các mệnh đề sau:}{%
  \y{Gọi $M$, $N$ là trung điểm của $AB$, $AC$ thì $\overrightarrow{MN}=(1;2;\dfrac{5}{2})$.}{d}
  \y{Điểm $M$ thoả mãn $\overrightarrow{MA}=3\overrightarrow{MB}$ có tọa độ $(0;0;-2)$.}{d}
  \y{Điểm $M$ thoả mãn $\overrightarrow{MA}=-2\overrightarrow{MB}$ có tọa độ $(-\dfrac{5}{3};0;\dfrac{2}{3})$.}{s}
  \y{Tâm của hình bình hành $ABCD$ có tọa độ $(-1;1;\dfrac{5}{2})$.}{s}}

\cauDS[id=C2-B7-D06-VD-038,muc=VD,loigiai={Kiểm tra từng mệnh đề bằng công thức tọa độ. a) Đúng. b) Sai. c) Đúng. d) Sai.}]{Trong không gian $Oxyz$, cho $A(4;4;3)$, $B(4;6;6)$, $C(2;4;-2)$, $D(-3;-4;6)$. Xét tính đúng sai của các mệnh đề sau:}{%
  \y{Gọi $M$, $N$ là trung điểm của $AB$, $AC$ thì $\overrightarrow{MN}=(-1;-1;-4)$.}{d}
  \y{Gọi $M$, $N$ là trung điểm của $AB$, $AC$ thì $\overrightarrow{MN}=(-1;-1;4)$.}{s}
  \y{Trọng tâm tứ diện $ABCD$ có tọa độ $(\dfrac{7}{4};\dfrac{5}{2};\dfrac{13}{4})$.}{d}
  \y{Điểm $M$ thoả mãn $\overrightarrow{MA}=\dfrac{1}{2}\overrightarrow{MB}$ có tọa độ $(4;-2;0)$.}{s}}

\cauDS[id=C2-B7-D06-VD-039,muc=VD,loigiai={Kiểm tra từng mệnh đề bằng công thức tọa độ. a) Sai. b) Sai. c) Đúng. d) Đúng.}]{Trong không gian $Oxyz$, cho $A(-1;2;-2)$, $B(1;-6;-4)$, $C(0;-1;3)$, $D(4;1;-5)$. Xét tính đúng sai của các mệnh đề sau:}{%
  \y{Tâm của hình bình hành $ABCD$ có tọa độ $(\dfrac{1}{2};\dfrac{1}{2};-\dfrac{1}{2})$.}{s}
  \y{Trọng tâm tứ diện $ABCD$ có tọa độ $(1;1;-2)$.}{s}
  \y{Tâm của hình bình hành $ABCD$ có tọa độ $(-\dfrac{1}{2};\dfrac{1}{2};\dfrac{1}{2})$.}{d}
  \y{Gọi $M$, $N$ là trung điểm của $AB$, $AC$ thì $\overrightarrow{MN}=(-\dfrac{1}{2};\dfrac{5}{2};\dfrac{7}{2})$.}{d}}

\cauDS[id=C2-B7-D06-VD-040,muc=VD,loigiai={Kiểm tra từng mệnh đề bằng công thức tọa độ. a) Sai. b) Sai. c) Đúng. d) Đúng.}]{Trong không gian $Oxyz$, cho $A(5;-5;-3)$, $B(1;-5;-5)$, $C(6;4;-2)$, $D(3;4;-4)$. Xét tính đúng sai của các mệnh đề sau:}{%
  \y{Gọi $M$, $N$ là trung điểm của $AB$, $AC$ thì $\overrightarrow{MN}=(-5;-9;-3)$.}{s}
  \y{Gọi $M$, $N$ là trung điểm của $AB$, $AC$ thì $\overrightarrow{MN}=(\dfrac{5}{2};-\dfrac{9}{2};\dfrac{3}{2})$.}{s}
  \y{Tâm của hình bình hành $ABCD$ có tọa độ $(\dfrac{11}{2};-\dfrac{1}{2};-\dfrac{5}{2})$.}{d}
  \y{Trọng tâm tứ diện $ABCD$ có tọa độ $(\dfrac{15}{4};-\dfrac{1}{2};-\dfrac{7}{2})$.}{d}}

\cauDS[id=C2-B7-D06-VD-041,muc=VD,loigiai={Kiểm tra từng mệnh đề bằng công thức tọa độ. a) Sai. b) Sai. c) Sai. d) Đúng.}]{Trong không gian $Oxyz$, cho $A(6;4;1)$, $B(-6;4;4)$, $C(0;4;1)$, $D(5;-3;-1)$. Xét tính đúng sai của các mệnh đề sau:}{%
  \y{Tâm của hình bình hành $ABCD$ có tọa độ $(1;4;3)$.}{s}
  \y{Điểm $M$ thoả mãn $\overrightarrow{MA}=-1\overrightarrow{MB}$ có tọa độ $(0;-4;\dfrac{5}{2})$.}{s}
  \y{Trọng tâm tứ diện $ABCD$ có tọa độ $(\dfrac{5}{4};\dfrac{9}{4};\dfrac{9}{4})$.}{s}
  \y{Tâm của hình bình hành $ABCD$ có tọa độ $(3;4;1)$.}{d}}

\cauDS[id=C2-B7-D06-VD-042,muc=VD,loigiai={Kiểm tra từng mệnh đề bằng công thức tọa độ. a) Sai. b) Sai. c) Đúng. d) Đúng.}]{Trong không gian $Oxyz$, cho $A(-1;1;6)$, $B(3;-2;4)$, $C(-2;5;2)$. Xét tính đúng sai của các mệnh đề sau:}{%
  \y{Tâm của hình bình hành $ABCD$ có tọa độ $(-\dfrac{3}{2};2;4)$.}{s}
  \y{Điểm $M$ thoả mãn $\overrightarrow{MA}=-2\overrightarrow{MB}$ có tọa độ $(\dfrac{5}{3};-1;\dfrac{11}{3})$.}{s}
  \y{Gọi $M$, $N$ là trung điểm của $AB$, $AC$ thì $\overrightarrow{MN}=(-\dfrac{5}{2};\dfrac{7}{2};-1)$.}{d}
  \y{Điểm $M$ thoả mãn $\overrightarrow{MA}=\dfrac{1}{2}\overrightarrow{MB}$ có tọa độ $(-5;4;8)$.}{d}}

\cauDS[id=C2-B7-D06-VD-043,muc=VD,loigiai={Kiểm tra từng mệnh đề bằng công thức tọa độ. a) Sai. b) Đúng. c) Sai. d) Sai.}]{Trong không gian $Oxyz$, cho $A(-3;-1;2)$, $B(-5;1;3)$, $C(-2;-3;1)$, $D(-5;0;3)$. Xét tính đúng sai của các mệnh đề sau:}{%
  \y{Tâm của hình bình hành $ABCD$ có tọa độ $(\dfrac{5}{2};2;-\dfrac{3}{2})$.}{s}
  \y{Gọi $M$, $N$ là trung điểm của $AB$, $AC$ thì $\overrightarrow{MN}=(\dfrac{3}{2};-2;-1)$.}{d}
  \y{Trọng tâm tứ diện $ABCD$ có tọa độ $(-\dfrac{15}{4};-\dfrac{3}{4};\dfrac{1}{4})$.}{s}
  \y{Trọng tâm tứ diện $ABCD$ có tọa độ $(-\dfrac{19}{4};-\dfrac{3}{4};\dfrac{9}{4})$.}{s}}

\cauDS[id=C2-B7-D06-VD-044,muc=VD,loigiai={Kiểm tra từng mệnh đề bằng công thức tọa độ. a) Đúng. b) Sai. c) Đúng. d) Sai.}]{Trong không gian $Oxyz$, cho $A(2;1;4)$, $B(-4;-2;2)$, $C(1;0;5)$, $D(-2;-2;-1)$. Xét tính đúng sai của các mệnh đề sau:}{%
  \y{Trọng tâm tứ diện $ABCD$ có tọa độ $(-\dfrac{3}{4};-\dfrac{3}{4};\dfrac{5}{2})$.}{d}
  \y{Tâm của hình bình hành $ABCD$ có tọa độ $(\dfrac{3}{2};\dfrac{1}{2};\dfrac{7}{2})$.}{s}
  \y{Điểm $M$ thoả mãn $\overrightarrow{MA}=3\overrightarrow{MB}$ có tọa độ $(-7;-\dfrac{7}{2};1)$.}{d}
  \y{Tâm của hình bình hành $ABCD$ có tọa độ $(\dfrac{9}{2};\dfrac{1}{2};\dfrac{3}{2})$.}{s}}

\cauDS[id=C2-B7-D06-VD-045,muc=VD,loigiai={Kiểm tra từng mệnh đề bằng công thức tọa độ. a) Đúng. b) Sai. c) Sai. d) Sai.}]{Trong không gian $Oxyz$, cho $A(4;5;-4)$, $B(-5;4;4)$, $C(2;6;1)$, $D(4;5;-2)$. Xét tính đúng sai của các mệnh đề sau:}{%
  \y{Trọng tâm tứ diện $ABCD$ có tọa độ $(\dfrac{5}{4};5;-\dfrac{1}{4})$.}{d}
  \y{Điểm $M$ thoả mãn $\overrightarrow{MA}=-\dfrac{1}{2}\overrightarrow{MB}$ có tọa độ $(1;\dfrac{17}{3};-\dfrac{4}{3})$.}{s}
  \y{Tâm của hình bình hành $ABCD$ có tọa độ $(-3;-\dfrac{11}{2};\dfrac{3}{2})$.}{s}
  \y{Điểm $M$ thoả mãn $\overrightarrow{MA}=\dfrac{1}{2}\overrightarrow{MB}$ có tọa độ $(13;6;-14)$.}{s}}

\cauDS[id=C2-B7-D06-VD-046,muc=VD,loigiai={Kiểm tra từng mệnh đề bằng công thức tọa độ. a) Sai. b) Sai. c) Đúng. d) Đúng.}]{Trong không gian $Oxyz$, cho $A(-1;2;-5)$, $B(-6;3;4)$, $C(1;6;2)$, $D(-2;-1;-1)$. Xét tính đúng sai của các mệnh đề sau:}{%
  \y{Điểm $M$ thoả mãn $\overrightarrow{MA}=-\dfrac{1}{2}\overrightarrow{MB}$ có tọa độ $(-\dfrac{8}{3};\dfrac{1}{3};-2)$.}{s}
  \y{Trọng tâm tứ diện $ABCD$ có tọa độ $(-2;-\dfrac{5}{2};0)$.}{s}
  \y{Gọi $M$, $N$ là trung điểm của $AB$, $AC$ thì $\overrightarrow{MN}=(\dfrac{7}{2};\dfrac{3}{2};-1)$.}{d}
  \y{Trọng tâm tứ diện $ABCD$ có tọa độ $(-2;\dfrac{5}{2};0)$.}{d}}

\cauDS[id=C2-B7-D06-VD-047,muc=VD,loigiai={Kiểm tra từng mệnh đề bằng công thức tọa độ. a) Sai. b) Sai. c) Sai. d) Đúng.}]{Trong không gian $Oxyz$, cho $A(1;2;5)$, $B(-5;0;-6)$, $C(-4;5;-5)$, $D(-2;-1;1)$. Xét tính đúng sai của các mệnh đề sau:}{%
  \y{Điểm $M$ thoả mãn $\overrightarrow{MA}=-1\overrightarrow{MB}$ có tọa độ $(-2;-1;-\dfrac{1}{2})$.}{s}
  \y{Trọng tâm tứ diện $ABCD$ có tọa độ $(-\dfrac{5}{2};-\dfrac{5}{4};\dfrac{3}{2})$.}{s}
  \y{Tâm của hình bình hành $ABCD$ có tọa độ $(-\dfrac{3}{2};0;\dfrac{7}{2})$.}{s}
  \y{Tâm của hình bình hành $ABCD$ có tọa độ $(-\dfrac{3}{2};\dfrac{7}{2};0)$.}{d}}

\cauDS[id=C2-B7-D06-VD-048,muc=VD,loigiai={Kiểm tra từng mệnh đề bằng công thức tọa độ. a) Đúng. b) Sai. c) Đúng. d) Sai.}]{Trong không gian $Oxyz$, cho $A(-1;-6;-3)$, $B(-4;1;-4)$, $C(-6;-5;-2)$, $D(6;4;0)$. Xét tính đúng sai của các mệnh đề sau:}{%
  \y{Điểm $M$ thoả mãn $\overrightarrow{MA}=-\dfrac{1}{2}\overrightarrow{MB}$ có tọa độ $(-2;-\dfrac{11}{3};-\dfrac{10}{3})$.}{d}
  \y{Gọi $M$, $N$ là trung điểm của $AB$, $AC$ thì $\overrightarrow{MN}=(-1;3;1)$.}{s}
  \y{Gọi $M$, $N$ là trung điểm của $AB$, $AC$ thì $\overrightarrow{MN}=(-1;-3;1)$.}{d}
  \y{Trọng tâm tứ diện $ABCD$ có tọa độ $(\dfrac{5}{4};\dfrac{3}{2};\dfrac{9}{4})$.}{s}}

\cauTLN[id=C2-B7-D06-VD-049,muc=VD,loigiai={Vì $M,N,P$ là trung điểm $BC,CA,AB$ nên $APMN$ là hình bình hành: $\overrightarrow{AP}=\overrightarrow{NM}$ hay $A=N+P-M=(-3;-6;-4)$; tương tự $B=M+P-N=(3;4;0)$, $C=M+N-P=(3;-6;-2)$.}]{Tam giác $ABC$ có các trung điểm của $BC$, $CA$, $AB$ lần lượt là $M(3;-1;-1)$, $N(0;-6;-3)$, $P(0;-1;-2)$. Tính $x_C+y_C+z_C$.}{-5}

\cauTLN[id=C2-B7-D06-VD-050,muc=VD,loigiai={Vì $M,N,P$ là trung điểm $BC,CA,AB$ nên $APMN$ là hình bình hành: $\overrightarrow{AP}=\overrightarrow{NM}$ hay $A=N+P-M=(4;5;2)$; tương tự $B=M+P-N=(4;-1;4)$, $C=M+N-P=(4;3;2)$.}]{Tam giác $ABC$ có các trung điểm của $BC$, $CA$, $AB$ lần lượt là $M(4;1;3)$, $N(4;4;2)$, $P(4;2;3)$. Tính $x_B+y_B+z_B$.}{7}

\cauTLN[id=C2-B7-D06-VD-051,muc=VD,loigiai={Vì $M,N,P$ là trung điểm $BC,CA,AB$ nên $APMN$ là hình bình hành: $\overrightarrow{AP}=\overrightarrow{NM}$ hay $A=N+P-M=(2;0;-6)$; tương tự $B=M+P-N=(-2;0;-4)$, $C=M+N-P=(-2;2;-2)$.}]{Tam giác $ABC$ có các trung điểm của $BC$, $CA$, $AB$ lần lượt là $M(-2;1;-3)$, $N(0;1;-4)$, $P(0;0;-5)$. Tính $x_B+y_B+z_B$.}{-6}

\cauTLN[id=C2-B7-D06-VD-052,muc=VD,loigiai={Vì $M,N,P$ là trung điểm $BC,CA,AB$ nên $APMN$ là hình bình hành: $\overrightarrow{AP}=\overrightarrow{NM}$ hay $A=N+P-M=(4;-5;0)$; tương tự $B=M+P-N=(0;-5;-6)$, $C=M+N-P=(2;3;-2)$.}]{Tam giác $ABC$ có các trung điểm của $BC$, $CA$, $AB$ lần lượt là $M(1;-1;-4)$, $N(3;-1;-1)$, $P(2;-5;-3)$. Tính $x_B+y_B+z_B$.}{-11}

\cauTLN[id=C2-B7-D06-VD-053,muc=VD,loigiai={Tâm $I$ là trung điểm của $AC$ và $BD$ nên $D=2I-B=(-5;-4;3)$.}]{Hình bình hành $ABCD$ có tâm $I(-1;-\dfrac{7}{2};2)$ và hai đỉnh $A(-4;-3;1)$, $B(3;-3;1)$. Tính $x_D+y_D+z_D$.}{-6}

\cauTLN[id=C2-B7-D06-VD-054,muc=VD,loigiai={Tâm $I$ là trung điểm của $AC$ và $BD$ nên $D=2I-B=(-4;1;-2)$.}]{Hình bình hành $ABCD$ có tâm $I(-\dfrac{3}{2};-\dfrac{3}{2};-\dfrac{1}{2})$ và hai đỉnh $A(0;-3;3)$, $B(1;-4;1)$. Tính $x_D+y_D+z_D$.}{-5}

\cauTLN[id=C2-B7-D06-VD-055,muc=VD,loigiai={Tâm $I$ là trung điểm của $AC$ và $BD$ nên $D=2I-B=(-7;-2;-5)$.}]{Hình bình hành $ABCD$ có tâm $I(-\dfrac{5}{2};-2;-1)$ và hai đỉnh $A(-4;-2;2)$, $B(2;-2;3)$. Tính $x_D+y_D+z_D$.}{-14}

\cauTLN[id=C2-B7-D06-VD-056,muc=VD,loigiai={Tâm $I$ là trung điểm của $AC$ và $BD$ nên $D=2I-B=(3;0;-1)$.}]{Hình bình hành $ABCD$ có tâm $I(1;2;-\dfrac{5}{2})$ và hai đỉnh $A(-2;3;-1)$, $B(-1;4;-4)$. Tính $x_D+y_D+z_D$.}{2}

\cauTLN[id=C2-B7-D06-VD-057,muc=VD,loigiai={Vì $G$ là trọng tâm tứ diện nên $A+B+C+D=4G$, suy ra $D=4G-A-B-C=(-1;-1;2)$.}]{Tứ diện $ABCD$ có $A(3;1;3)$, $B(-1;-4;-3)$, $C(1;3;1)$ và trọng tâm $G(\dfrac{1}{2};-\dfrac{1}{4};\dfrac{3}{4})$. Tính $x_D+y_D+z_D$.}{0}

\cauTLN[id=C2-B7-D06-VD-058,muc=VD,loigiai={Vì $G$ là trọng tâm tứ diện nên $A+B+C+D=4G$, suy ra $D=4G-A-B-C=(4;3;-4)$.}]{Tứ diện $ABCD$ có $A(4;-4;3)$, $B(4;3;3)$, $C(4;-1;2)$ và trọng tâm $G(4;\dfrac{1}{4};1)$. Tính $x_D+y_D+z_D$.}{3}

\cauTLN[id=C2-B7-D06-VD-059,muc=VD,loigiai={Vì $G$ là trọng tâm tứ diện nên $A+B+C+D=4G$, suy ra $D=4G-A-B-C=(2;-4;3)$.}]{Tứ diện $ABCD$ có $A(-3;-3;4)$, $B(1;0;-3)$, $C(1;1;0)$ và trọng tâm $G(\dfrac{1}{4};-\dfrac{3}{2};1)$. Tính $x_D+y_D+z_D$.}{1}

\cauTLN[id=C2-B7-D06-VD-060,muc=VD,loigiai={Vì $G$ là trọng tâm tứ diện nên $A+B+C+D=4G$, suy ra $D=4G-A-B-C=(-4;-2;-1)$.}]{Tứ diện $ABCD$ có $A(2;4;3)$, $B(-3;-2;2)$, $C(3;0;-1)$ và trọng tâm $G(-\dfrac{1}{2};0;\dfrac{3}{4})$. Tính $x_D+y_D+z_D$.}{-7}
